% Modernised reading — fonds Alexandre Grothendieck, shelfmark 161-4
% (pages 1-31, the whole folder).
%
% This is an INTERPRETATION, not a transcription. It re-reads the pages in
% today's notation and vocabulary, states the hypotheses the manuscript leaves
% implicit, and drops the critical apparatus so that the mathematics reads
% straight through. Everything the brackets and underlines carried moves into a
% footnote. It is held to being mathematically correct as it stands; where the
% manuscript is loose or mistaken, the true statement is what appears here and
% the footnote says what the page has. In any dispute of substance the
% transcription governs: whoever cites Grothendieck must cite
% batch-01.fr.tex (pp. 1-20) or batch-02.fr.tex (pp. 21-31).
%
% Pass: Opus 5.5 (claude-opus-5-5), 2026-09-22 — first pass, unchecked against the pages by a human.
% The folder was read whole, both batches together; this is one document for
% the shelfmark. It was made on Opus 5.5 and has not been measured against the
% reference reading of folder 115 (made on Opus 5).
%
% Spine. Seven strands, not one argument.
%   I    pp. 2-6 (his pp. 1-5): faithfully flat descent, items (1)-(5); (5)
%        is a heading with nothing under it.
%   II   pp. 7-9: radicial / universally open morphisms, quasi-affine
%        morphisms and the boxed implication table, ample sheaves (EGA II 4.5).
%   III  pp. 10-12: course plans — stability rules a) b) c) for a class M,
%        syllabi, the pencil plan (d)-(h) with P(M) = V(M)^*/G_m.
%   IV   pp. 14, 16: schemes as functors; Hom(X, Spec A); the sites diagram;
%        representability by Zariski sheaves covered by open affines.
%   V    pp. 22-24, 26-28: the course outline, English prerequisites then
%        items (1)-(10) on affine k-schemes as functors V_A.
%   VI   pp. 28-31 and 21: item (11), the spectrum; p. 21 by content after
%        the outline (the « 10') Fonctorialité » of p. 30's margin).
%   VII  pp. 19-20: « Compléments » on limits and sums in Aff_k; they use the
%        outline's e_k, V_A and the spectra of (11), so they are read last.
%   VIII pp. 17, 15, 25: not the course. Product of topoi (p. 15 continues
%        p. 17); p. 25 pencil fragments. Cross-reference to 161-6 p. 22.
%
% Notation fixed for the whole document. k a commutative ring (often a
% field). V_A = Spec A seen as the functor k' -> Hom_{k-alg}(A, k'); E^I_k his
% affine space functor k' -> k'^I; e_k = V_k, emptyset_k = V_0; Aff_k, Alg_k,
% Sch_k. His « espace algébrique affine » = affine k-scheme (NOT Artin's
% algebraic spaces). Descent: f : S' -> S, g : X -> Y an S-morphism, primes =
% base change to S', S'' = S' x_S S'. « cohérent » (his) = qcqs. A class of
% morphisms M with his rules a) b) c). Topos pages: pi(C, E) = Fun(C^op, E),
% hat = presheaves, tilde = sheaves. On A, f < g (his ≺) iff f^n in gA.
%
% Substantive departures, with pages.
%  p. 2   lying-over from A/p -> B/pB injective: the step needs B (x) k(p) != 0,
%         i.e. faithful flatness applied to k(p); supplied.
%  p. 2   margin criterion « M fid. plat <=> M (x) k(p) != 0 »: M flat supplied.
%  p. 3   flat base change lemma: f qcqs and F quasi-coherent supplied.
%  p. 4   f, f' for g, g' (flagged in the transcription); quasi-compactness
%         descends along f surjective qc, the rest along f fpqc (stated so).
%  p. 4   descent of immersions stated for open, closed, quasi-compact
%         immersions only; arbitrary immersions left as the page's own open
%         question (pp. 3, 5).
%  p. 5   quasi-affine: « X -> Spec g_* O_X immersion » stated as open
%         immersion (equivalent here); Ssf(F) read as quasi-coherent submodules.
%  p. 6   ampleness criterion written with g : X -> Y, not f, F on X not S.
%  p. 6   the NB on projective / quasi-projective: « among the listed
%         properties »; reason given (the sheaf, not its ampleness, fails to
%         descend).
%  p. 7   conditions a)-d) renumbered (i)-(iv); a) b) c) kept for p. 10's rules.
%  p. 8   projective = proper + quasi-projective: Y quasi-compact supplied;
%         definitions are EGA II's, footnoted; finite = proper quasi-finite
%         holds without noetherian hypothesis.
%  p. 9   composite of ample sheaves: Z quasi-compact supplied (uniform n_0).
%  p. 10  M contains isomorphisms supplied; the last equivalence quantified
%         over base changes of g; « qu.-coh. » for « cohérent » (transcription).
%  p. 11  identity principle: the source Z is reduced, not X.
%  p. 12  « de type fini » includes quasi-separated, unlike EGA; flagged.
%  p. 14  Γ(U_f, O_X) read Γ(φ^{-1}(U_f), O_X); equivalence of Zariski-sheaf
%         categories Sch_k vs Aff_k supplied (page draws the arrow only).
%  p. 16  Cor: « soit un faisceau » for « soit un schéma » (flagged sic).
%  p. 17  the characterisation of Ĉ x D~ by p_*, q_* alone is false
%         (counterexample given); corrected to all partial functors; p. 15 takes
%         the correct route. γ(C,E) is in fact always an equivalence (classical).
%  p. 19  his k[I] (function algebra) written k^I; k[I] today = group algebra.
%  p. 20  Ker φ -> ker u (transcription); induced topology stated only for the
%         cases the page names.
%  p. 21  Spec ∏ A_i = ⊔ Spec A_i for I finite only; « [vrai sur un corps alg.
%         clos] » read as the maximal spectrum, as the next line says;
%         φ^{-1}(p) read u^{-1}(p).
%  p. 23  « sur A » for « sur k » (transcription).
%  p. 27  idempotent description of Hom(∏ A_i, k') for I finite only.
%  p. 28  « O.K. si R fini »: modern effectivity needs R finite locally free.
%  p. 29  « contiennent A » for « S » (transcription).
%  p. 30  « f_i ∈ I » (transcription); finite subfamily renamed I_0.
%
% Announced and never established on the pages: (5) « Questions d'effectivité
% de données de descente » (p. 6) is a bare heading; « Expliquer ce que cela
% signifie » (p. 5); the converse of p. 7's Prop (« nettement plus difficile »);
% item (7) of p. 24 (illegible); the passage to the quotient promised p. 28;
% the sheaf property of p. 31 (« … »); on the topos pages, that γ is an
% equivalence (p. 17 ends « ce qui prouvera… », p. 15 reduces to two cases and
% stops). Dating: no leaf carries a date. The one outside reference, « EGA II
% 4.5 » in the margin of p. 9, places that sheet (of the pp. 7-9 run) after
% EGA II and dates nothing else — not the descent pages, the course pages or
% the topos leaves. P. 22 says only that the course was a two-part summer
% course taught in English; no course or year is named.
%
% Revised 2026-09-23 (Opus 5.5): header dating narrowed to the p. 9 sheet;
% résumé no longer calls the descent stratum « plus ancienne peut-être »
% (nothing on the leaves orders it), and « Le fil du dossier » says the same;
% section VIII states the link with 161-6 p. 22 without an order (« reprend
% exactement là » removed) and footnotes the c) lettering mismatch; p. 14
% footnote records that « fid. » is read uncertain; p. 16 footnote now says the
% transcription flags the slip (sic note added there); section VII footnotes the
% page's gothic/roman letters on p. 20 (batch 1 corrected to match).
\documentclass[11pt,a4paper]{article}
% Le préambule commun à toutes les transcriptions du fonds Grothendieck.
%
% Il tient en une idée : **l'incertitude se compose, elle ne se résout pas.**
% Une transcription qui lisse un mot douteux a détruit la seule chose pour
% laquelle on la faisait — un lecteur doit pouvoir savoir, à chaque endroit,
% ce qui était sur le feuillet et ce qui a été ajouté. Les six macros qui
% suivent sont donc l'appareil critique, et non de la décoration.
%
% Les mêmes commandes sont reconnues par `scripts/render.mjs`, qui produit la
% vue de lecture du site. Une macro ajoutée ici doit l'être aussi là-bas,
% faute de quoi le rendu échouera bruyamment — ce qui est voulu : un rendu
% silencieusement faux est pire qu'un rendu absent.

\ProvidesPackage{grothendieck}[2026/08/08 Transcriptions du fonds Grothendieck]

\usepackage{amsmath,amssymb,amsthm}
\usepackage{mathrsfs}
\usepackage{stmaryrd}
% Les diagrammes commutatifs sont partout dans ce fonds : c'est la langue
% même de la géométrie algébrique de Grothendieck, et une transcription qui
% les rendrait en prose ne transcrirait rien.
\usepackage{tikz-cd}
% Français par défaut — c'est la langue des feuillets — mais l'anglais est
% chargé aussi : la traduction et le résumé sont des documents anglais, et la
% typographie française (espaces avant les deux-points) y serait une faute.
% Ces documents basculent par \selectlanguage{english} en tête de corps.
\usepackage[english,french]{babel}
\usepackage{fontspec}
\usepackage{geometry}
\geometry{a4paper,margin=25mm,marginparwidth=28mm}
\usepackage{marginnote}
\usepackage{xcolor}
% ulem, not soul. `soul`'s \st and \ul are built on a character-by-character
% scan that XeLaTeX's font machinery breaks: the document fails with an
% undefined control sequence rather than degrading. ulem does the same two jobs
% and is Unicode-engine safe. `normalem` stops it hijacking \emph, which the
% transcriptions use for Grothendieck's own emphasis.
\usepackage[normalem]{ulem}
\usepackage{hyperref}
\hypersetup{colorlinks=true,linkcolor=black,urlcolor=black,pdfborder={0 0 0}}

% --- Métadonnées du lot -----------------------------------------------------
% Reprises telles quelles par le script de rendu, qui les affiche en tête de
% la vue de lecture. Elles fixent ce que le fichier prétend couvrir : sans
% elles, une transcription flotte sans point d'attache dans un fonds de seize
% mille feuillets.

\newcommand{\folder}[1]{\gdef\grothFolder{#1}}
\newcommand{\batch}[1]{\gdef\grothBatch{#1}}
\newcommand{\leaves}[2]{\gdef\grothFirst{#1}\gdef\grothLast{#2}}
\newcommand{\foldertitle}[1]{\gdef\grothFoldertitle{#1}}
\gdef\grothFolder{?}\gdef\grothBatch{?}\gdef\grothFirst{?}\gdef\grothLast{?}\gdef\grothFoldertitle{}

% --- Le résumé --------------------------------------------------------------
%
% Il ouvre la lecture modernisée, sous le titre « Résumé », et il est composé
% en petit. Ce n'est pas un ornement : c'est le seul passage du document qui
% ne soit pas des mathématiques, et un lecteur doit pouvoir le reconnaître
% avant de l'avoir lu — puis le sauter s'il connaît déjà le sujet, ou s'y
% arrêter s'il ne le connaît pas. Le filet qui le suit marque la reprise du
% corps.

\newenvironment{resume}
  {\small\setlength{\parskip}{0.45em}}
  {\par\medskip\hrule\medskip}

% --- Mots-clés ---------------------------------------------------------------
%
% En anglais, en clôture du résumé de la lecture modernisée : le vocabulaire
% moderne sous lequel chercher ce que les feuillets construisent. Le manifeste
% du site (`npm run manifest`) les extrait pour étiqueter la cote entière —
% cette ligne est la source unique des tags, il n'y a pas de fichier à côté.
\newcommand{\keywords}[1]{\par\smallskip\noindent
  {\footnotesize\textsc{Keywords} --- \textit{#1}}\par}

% --- L'appareil critique ----------------------------------------------------

% Le numéro de feuillet, en marge. C'est l'ancre qui permet de régler un
% désaccord sur une formule en regardant *un* feuillet plutôt qu'une chemise
% de six cents. Le site s'en sert aussi pour tourner les pages du fac-similé
% quand on fait défiler la transcription.
\newcommand{\leaf}[1]{%
  \marginnote{\footnotesize\color{gray!70}\textsf{#1}}%
  \label{leaf:#1}\ignorespaces}

% Illisible : on ne devine pas. Un mot inventé qui se lit comme les autres est
% le pire résultat possible.
\newcommand{\ill}{\textcolor{red!60!black}{[\,\ldots\,]}}

% Lecture douteuse : le mot est donné, et signalé comme incertain.
\newcommand{\uncertain}[1]{\uline{#1}}

% Ajout de l'éditeur — une lettre manquante, un mot sous-entendu.
\newcommand{\add}[1]{\textcolor{blue!50!black}{[#1]}}

% Biffé par Grothendieck lui-même. Ce qu'il a rayé fait partie du document :
% une rature dit souvent où la pensée a changé de direction.
\newcommand{\struck}[1]{\sout{#1}}

% Note du transcripteur, sur l'état matériel du feuillet.
\newcommand{\note}[1]{\par\smallskip{\small\color{gray!60!black}\textit{[#1]}}\par\smallskip}

% Note marginale de Grothendieck, distincte du corps du texte.
\newcommand{\marginal}[1]{\par\smallskip\noindent\hspace{1em}%
  {\small\color{gray!60!black}#1}\par\smallskip}

% --- Titre ------------------------------------------------------------------

\AtBeginDocument{%
  \begin{center}
    {\large\bfseries Fonds Alexandre Grothendieck --- cote n\textsuperscript{o}~\grothFolder}\\[2pt]
    {\itshape\grothFoldertitle}\\[4pt]
    {\small feuillets \grothFirst--\grothLast\ (lot \grothBatch)}\\[2pt]
    {\footnotesize\color{gray!60!black}Transcription établie d'après le fac-similé numérisé
     par l'Université de Montpellier.\\
     Les lectures incertaines sont soulignées, les illisibles notées [\,\ldots\,],
     les ajouts entre crochets.}
  \end{center}
  \vspace{1em}}

\endinput


\folder{161-4}
\pages{1}{31}
\dating{[après 1961] — le groupe « Dossiers rassemblés par Grothendieck sur différentes thématiques » (161-1 à 162-6) est daté [après 1961-vers 1977]}
\watermark{Édition de démonstration\\interprétation personnelle de l'œuvre}
\foldertitle{Introduction à géométrie algébrique : notes manuscrites (s.d.) — lecture modernisée du dossier entier}

\begin{document}

\section*{Résumé}

\begin{resume}
La chemise porte au crayon « Introduction à Géom. Alg. », et c'est bien ce
qu'elle contient : les papiers de travail de quelqu'un qui prépare un cours
d'initiation à la géométrie algébrique, et qui se demande par où commencer.
Les feuillets ne forment pas un texte suivi. Ce sont des morceaux de cours,
des plans raturés et refaits, des listes de ce qu'il ne faut pas oublier ; et,
glissés parmi eux, trois feuillets sur un tout autre sujet.

La géométrie algébrique étudie les ensembles de solutions des systèmes
d'équations polynomiales. La difficulté, pour qui l'enseigne, est de dire ce
qu'est l'« espace » formé par ces solutions. Il y a deux chemins. Le premier,
celui des \emph{Éléments de géométrie algébrique}, fabrique un espace
topologique muni d'un faisceau d'anneaux, puis établit une à une les propriétés
des morphismes entre de tels objets : il est complet, mais long, et un plan du
dossier en appelle la première étape, sans ironie apparente, la « traversée du
désert ». Le second chemin part d'une idée plus simple. On connaît une serrure
par l'ensemble des clés qu'elle accepte ; de même, on connaît un système
d'équations par l'ensemble de ses solutions dans \emph{chaque} anneau --- les
nombres rationnels, les réels, les complexes, les entiers modulo $5$, et tous
les autres --- et par la façon dont ces ensembles se transforment quand on
change d'anneau. Un cercle n'a pas les mêmes points sur $\mathbb{Q}$ et sur
$\mathbb{C}$ ; la donnée de tous ces ensembles de points à la fois, c'est le
cercle, et l'on ne perd rien à le définir ainsi. Les feuillets du cours
choisissent ce second chemin. Ils en tirent une conclusion que l'une des pages
énonce avec une satisfaction visible : la notion de schéma s'en trouve
« explicitée sans recours à la notion d'espaces topologiques ». L'espace
topologique n'est pas pour autant abandonné : il revient à la fin, reconstruit,
sous le nom de \emph{spectre}, et ses points sont les « lieux » où tombent les
solutions.

Une page en anglais situe le cours. C'est un cours d'été en deux parties, une
introduction à la géométrie algébrique suivie d'une introduction aux groupes
algébriques, et la page note que la première est le prérequis de la seconde.
Le plan porte la trace de cette destination : il construit tôt les noyaux, le
groupe spécial linéaire comme noyau du déterminant, les stabilisateurs ; et il
montre pourquoi les schémas affines ne suffisent pas pour les quotients, sur un
exemple où le quotient calculé dans la catégorie affine est un point --- « ce
qui est idiot ! ».

Une autre strate, d'un autre niveau, et que rien sur les feuillets ne situe
dans le temps par rapport au cours, fait la liste de ce qu'un cours plus avancé devrait contenir. La plus longue série
porte sur la \emph{descente}. L'image est celle d'un atlas : on connaît un pays
par des cartes locales et par la manière dont elles se recouvrent. La descente
généralise l'atlas à des « recouvrements » beaucoup plus grossiers que des
ouverts --- les morphismes fidèlement plats --- et pose deux questions : quelles
propriétés d'un objet peut-on vérifier après être passé au recouvrement, et
quels objets du recouvrement proviennent d'un objet d'en bas ? À la première,
une longue liste répond « presque toutes », avec une exception notée à part :
être projectif. La seconde reste un titre sans rien dessous.

Ces pages ont un autre intérêt, qui ne demande aucune préparation : elles
montrent un auteur qui juge son propre sujet. La matière est « messy », dit la
page anglaise, et l'on ne pourra en donner qu'un survol. Tous les schémas des
« Anciens », note un plan, étaient de type fini et séparés sur un corps ; ceux
qu'on rencontre aujourd'hui sont presque tous quasi-séparés, et certains, parmi
les importants, ne sont pas quasi-compacts. Ailleurs, devant une construction
qui ne donne que ce qu'on savait déjà, une parenthèse : « nous fait une belle
jambe ! ». Le souci de l'ordre est constant : les plans sont renumérotés,
les rubriques déplacées par des flèches, une marge promet de voir « tôt » les
schémas non séparés.

Les trois feuillets étrangers au cours portent sur les topos. Deux d'entre eux
cherchent à montrer que les préfaisceaux sur une petite catégorie à valeurs
dans un topos forment le produit de deux topos ; ils s'arrêtent juste avant la
fin. La même question est posée, sur une page barrée, dans le dossier 161-6.

Les noms à chercher ensuite : descente fidèlement plate, foncteur des points,
faisceau de Zariski et critère de représentabilité, schémas en groupes affines,
espaces sobres et spectraux, produit de topos.

\keywords{faithfully flat descent, fpqc descent, effective descent, universally submersive morphism, radicial morphism, universal homeomorphism, quasi-affine morphism, relatively ample invertible sheaf, permanence properties of morphisms, projective bundle, functor of points, Zariski sheaf, representability criterion, affine group scheme, constant group scheme, Nullstellensatz, limits of affine schemes, Zariski topology, sober space, spectral space, structure sheaf, product of topoi, presheaf topos}
\end{resume}

\pagerange{1}{31}
\section*{Le fil du dossier, et les conventions}

\emph{Les strates.} Le dossier est un assemblage de feuillets de papiers, d'encres
et de moments différents, et l'ordre dans lequel les archivistes les ont
numérotés n'est pas toujours celui de l'argument. Après la chemise, il se lit
en sept ensembles de feuillets, qui sont les sept sections qui suivent.

\begin{enumerate}
  \item \emph{Page 1}, la chemise : « Introduction à Géom. Alg. », au crayon,
  sur l'onglet.
  \item \emph{Pages 2 à 6} (paginées par lui 1 à 5, encre bleue) : la descente
  fidèlement plate, en rubriques numérotées (1) à (5). La rubrique (5) n'a que
  son titre. Section I.
  \item \emph{Pages 7 à 9}, feuilles volantes de même main : morphismes
  radiciels et universellement ouverts, morphismes quasi-affines avec un tableau
  encadré d'implications, faisceaux amples. Section II.
  \item \emph{Pages 10 à 12} : trois plans de cours, très surchargés. Les
  règles de stabilité a), b), c) d'une classe de morphismes, deux programmes
  numérotés, puis au crayon un plan à rubriques (d) à (h) où apparaissent les
  schémas comme foncteurs et $P(M)$. Section III.
  \item \emph{Pages 14 et 16} : les schémas comme foncteurs --- morphismes
  d'un espace localement annelé vers un schéma affine, le diagramme des
  catégories de faisceaux, le critère de représentabilité. Section IV.
  \item \emph{Pages 22 à 24 et 26 à 28} : le cours lui-même. Une page de
  prérequis en anglais, puis un plan en français, rubriques (1) à (10), sur les
  schémas affines sur $k$ vus comme foncteurs sur les $k$-algèbres. Section V.
  \item \emph{Pages 28 à 31, et 21} : la rubrique (11), le spectre ---
  ensemble sous-jacent, topologie, faisceau structural. La page 21, que les
  archivistes placent avant le plan, emploie la notation du plan ($V_A$, les
  lettres gothiques) et traite la fonctorialité du spectre, qu'une marge de la
  page 30 annonce comme une rubrique « 10') » à insérer ; par son contenu elle
  vient donc après les pages 22 à 28. Section VI.
  \item \emph{Pages 19 et 20}, feuille titrée « Géom. Algébrique » :
  « Compléments sur le formalisme des $\varprojlim$ et des sommes dans
  $\mathrm{Aff}_k$ ». Elles emploient l'objet final $e_k$ et la notation $V_A$
  des pages 22 à 26, et parlent des spectres de la rubrique (11) : on les lit en
  dernier. Section VII.
  \item \emph{Pages 15, 17 et 25} : étrangères au cours. Les pages 15 et 17 sur
  le produit de deux topos --- la page 15, qui commence par « Or », continue la
  page 17 et non l'inverse ---, la page 25 au crayon, sur des topos aussi, lue
  seulement par fragments. Section VIII.
\end{enumerate}

Les pages 13 et 18 sont blanches. Rien sur les pages ne date le cours ni ne le
situe : on sait seulement qu'une marge de la page 9 renvoie à EGA II, ce qui
place cette feuille après EGA II sans rien dire des autres, et que la page 22
prévoit un cours d'été en deux parties, enseigné en anglais. On ne
cherche pas à l'identifier.

\emph{Deux niveaux.} Les sections I à III supposent connue la théorie des
schémas et en organisent les notions ; les sections IV à VII construisent
l'objet depuis le début. Rien n'indique que les deux strates appartiennent au
même cours, et on ne les confond pas.

\emph{Notations.} Elles valent pour tout le document.

\begin{itemize}
  \item $k$ est un anneau commutatif, souvent un corps ; $\mathrm{Alg}_k$ la
  catégorie des $k$-algèbres commutatives, $\mathrm{Aff}_k$ celle des schémas
  affines sur $k$, $\mathrm{Sch}_k$ celle des $k$-schémas.
  \item Pour une $k$-algèbre $A$, $V_A$ est le foncteur
  $k' \mapsto \operatorname{Hom}_{k\text{-alg}}(A, k')$ sur $\mathrm{Alg}_k$ ;
  c'est $\operatorname{Spec} A$ vu par ses points. $\mathbb{E}^I_k$ est le
  foncteur $k' \mapsto k'^I$, c'est-à-dire l'espace affine $V_{P_I}$ avec
  $P_I = k[(T_i)_{i \in I}]$. On pose $e_k = V_k$ (objet final) et
  $\emptyset_k = V_0$.
  \item Ce que les pages appellent \emph{espace algébrique affine} sur $k$ est,
  dans la langue d'aujourd'hui, un $k$-schéma affine présenté par son foncteur
  des points.\footnote{Le mot ne doit pas être confondu avec les \emph{espaces
  algébriques} au sens d'Artin, introduits à la fin des années 1960, qui sont
  une généralisation des schémas. Ici l'adjectif « affine » ne quitte jamais le
  nom, et l'objet est exactement un foncteur représentable $V_A$.}
  \item Descente : $f : S' \to S$ est le morphisme de descente,
  $g : X \to Y$ un $S$-morphisme, un accent désigne le changement de base à
  $S'$ ($X' = X \times_S S'$, $g' = g \times_S S'$), et
  $S'' = S' \times_S S'$.
  \item Il dit \emph{cohérent} pour « quasi-compact et
  quasi-séparé » ;\footnote{C'est la terminologie de SGA 4 ; la page 10 la
  définit elle-même (« Morphismes cohérents $=$ qu.-cpt qu.-sép. »). Le mot a
  d'autres sens aujourd'hui (faisceau cohérent), d'où l'abréviation
  « qcqs » qu'on lui préfère.} on écrit « qcqs ». « fpqc » abrège
  « fidèlement plat et quasi-compact ».
  \item $\mathcal{M}$ désigne une classe de morphismes, et a), b), c) ses
  règles de stabilité, fixées à la section III.
  \item Topos : $\widehat{C}$ est le topos des préfaisceaux sur $C$,
  $\widetilde{D}$ celui des faisceaux sur $D$ pour une topologie donnée, et
  $\pi(C, E) = \underline{\operatorname{Hom}}(C^{\circ}, E)$ le topos des
  préfaisceaux sur $C$ à valeurs dans le topos $E$.
  \item Sur un anneau $A$, on écrit $f \prec g$ lorsque $f$ appartient à la
  racine de l'idéal $gA$ ; c'est la notation de la page 30.
\end{itemize}

\emph{Ce que le dossier annonce sans l'établir.} La rubrique (5) de la
descente, sur l'effectivité des données de descente, n'a que son titre. La
réciproque de la proposition de la page 7 est annoncée « nettement plus
difficile » et n'est pas donnée. Le passage au quotient est promis à la page 28.
Le caractère de faisceau du faisceau structural est laissé en points de
suspension à la page 31. Sur les pages de topos, l'énoncé visé --- que le
morphisme de comparaison $\gamma$ est une équivalence --- n'est pas atteint.

\pagerange{2}{6}
\section*{I. La descente fidèlement plate (pages 2 à 6)}

Cinq pages paginées de sa main, en rubriques numérotées : (1) descente de
propriétés de morphismes, (2) descente de propriétés topologiques, (3) descente
effective des faisceaux quasi-cohérents, des sous-schémas, des morphismes, des
schémas affines, (4) descente de l'amplitude, (5) questions d'effectivité.
Elles sont précédées de rappels sur la platitude. C'est, en abrégé, la matière
que l'on trouve aujourd'hui dans l'exposé VIII de SGA 1 et au paragraphe 2
d'EGA IV ; la page n'y renvoie pas.

\pagerange{2}{2}
\subsection*{Platitude fidèle (page 2)}

Soit $A \to B$ un morphisme d'anneaux. On dit que $B$ est \emph{fidèlement
plat} sur $A$ si $B$ est plat et si, pour tout $A$-module $M$,
$M \otimes_A B = 0$ entraîne $M = 0$.

\textbf{Proposition.} $B$ est fidèlement plat sur $A$ si et seulement si $B$
est plat et $\operatorname{Spec} B \to \operatorname{Spec} A$ est surjectif.

Si $B$ est fidèlement plat, alors pour tout $A$-module $M$ l'application
$M \to M \otimes_A B$ est injective ; en particulier $A \to B$ est injectif et
$\mathfrak{J} B \cap A = \mathfrak{J}$ pour tout idéal $\mathfrak{J}$ de $A$.
Pour un idéal premier $\mathfrak{p}$, la fibre $B \otimes_A k(\mathfrak{p})$
n'est pas nulle, puisque $k(\mathfrak{p}) \neq 0$ : il existe donc un idéal
premier de $B$ au-dessus de $\mathfrak{p}$.\footnote{La page conclut de
l'injectivité de $A/\mathfrak{p} \to B/\mathfrak{p}B$ à l'existence d'un
premier au-dessus de $\mathfrak{p}$. L'injectivité seule ne suffit pas :
$\mathbb{Z} \to \mathbb{Z}[1/2]$ est injectif et aucun premier de
$\mathbb{Z}[1/2]$ n'est au-dessus de $2$. Ce qui fait marcher l'argument est la
fidélité appliquée au module $k(\mathfrak{p})$, et c'est ce qu'on écrit. Un mot
cerclé, lu « Ind. » avec doute, précède la phrase sur la page.} Réciproquement,
si $B$ est plat et le morphisme des spectres surjectif, et si $M \neq 0$, on
choisit un sous-module monogène $A/I \neq 0$ de $M$ et un idéal maximal
$\mathfrak{m} \supset I$ ; par platitude $B/IB \subset M \otimes_A B$, et
$B/IB$ se surjecte sur $B/\mathfrak{m}B$, qui n'est pas nul puisque
$\mathfrak{m}$ est induit par un premier de $B$.\footnote{La page raisonne
ainsi, avec plusieurs reprises biffées (« on peut supposer $M$ de t.f. »), et la
réduction aux quotients $A/\mathfrak{m}$ y est lue avec doute.}

La marge de la page ajoute deux critères, sous une forme qu'on complète :
un $A$-module \emph{plat} $M$ est fidèlement plat si et seulement si
$M \otimes_A k(\mathfrak{p}) \neq 0$ pour tout premier $\mathfrak{p}$ ; et si
$A \to B$ est un homomorphisme local d'anneaux locaux, $B$ plat équivaut à
$B$ fidèlement plat.\footnote{La marge écrit le premier critère sans
l'hypothèse de platitude, que l'énoncé exige ; elle est dans le sens du
contexte, qui ne parle que de modules plats. La marge est écrite en travers de
la page, sur plusieurs colonnes, et d'autres bribes s'y lisent sans ordre sûr.}

Un morphisme de schémas $f : X \to Y$ est dit fidèlement plat s'il est plat et
surjectif.\footnote{La définition, sur la page, s'arrête après « ssi » ; on
donne la définition usuelle, qui est celle que la suite emploie.} Si $f$ est
plat, alors pour tout $x \in X$ d'image $y$, l'application
$\operatorname{Spec} \mathcal{O}_{X,x} \to \operatorname{Spec} \mathcal{O}_{Y,y}$
est surjective (un homomorphisme local plat est fidèlement plat) ; par
conséquent tout point maximal de $X$ --- point générique d'une composante
irréductible --- est au-dessus d'un point maximal de $Y$. C'est le théorème de
\emph{descente des générisations} : un morphisme plat est
\emph{générisant}.

\pagerange{3}{3}
\subsection*{La topologie de $S$, quotient de celle de $S'$ (page 3)}

\textbf{Corollaire.} Un morphisme plat et localement de présentation finie est
universellement ouvert.\footnote{« univ. » est ajouté au-dessus de la ligne.}

\textbf{Proposition.} Soit $f : S' \to S$ fidèlement plat. L'application
$T \mapsto f^{-1}(T)$, des sous-schémas de $S$ vers ceux de $S'$, est
injective.\footnote{La page ajoute une seconde propriété, illisible. L'injectivité
n'exige pas la quasi-compacité : les ensembles sous-jacents se séparent par
surjectivité, et les idéaux par platitude fidèle des anneaux locaux
$\mathcal{O}_{S,s} \to \mathcal{O}_{S',s'}$.} Par conséquent, pour des
$S$-schémas $X$ et $Y$, l'application
$\operatorname{Hom}_S(X, Y) \to \operatorname{Hom}_{S'}(X', Y')$ est injective
(on l'applique aux graphes).

\textbf{Lemme (changement de base plat).} Soit un carré cartésien de
morphismes $X' \to X$, $f : X \to Y$, $f' : X' \to Y'$, $Y' \to Y$, avec
$Y' \to Y$ plat. Si $f$ est quasi-compact et quasi-séparé et $F$ un
$\mathcal{O}_X$-module quasi-cohérent, l'homomorphisme canonique
$f_*(F)' \to f'_*(F')$ est un isomorphisme.\footnote{Les hypothèses « $f$
qcqs » et « $F$ quasi-cohérent » ne sont pas sur la page, qui écrit seulement le
carré et la flèche ; sans elles l'énoncé est faux.}

\textbf{Théorème.} Soit $f : S' \to S$ un morphisme plat, et $T \subset S$
l'image d'un morphisme quasi-compact $g : Z \to S$. Alors
\[
f^{-1}(\overline{T}) = \overline{f^{-1}(T)} .
\]

La page en donne deux démonstrations. La première se ramène au cas affine :
si $Z = \operatorname{Spec} C$, $S = \operatorname{Spec} A$ et
$u : A \to C$, l'adhérence $\overline{T}$ est $V(\operatorname{Ker} u)$, et la
platitude de $A \to A'$ donne
$\operatorname{Ker}(u \otimes A') = (\operatorname{Ker} u) A'$. La seconde
est plus éclairante, et la page note elle-même qu'elle « marche pour tout
morphisme générisant » : comme $g$ est quasi-compact, $\overline{T}$ est
l'ensemble des spécialisations des points de $T$, et de même pour
$f^{-1}(T)$, qui est l'image du morphisme quasi-compact
$Z \times_S S' \to S'$ ; si $s' \in f^{-1}(\overline{T})$, le point $f(s')$
est spécialisation d'un $t \in T$, et comme $f$ est générisant, $t$ se relève
en une générisation de $s'$, qui est dans $f^{-1}(T)$.\footnote{Les deux
ajouts de la colonne de droite qui portent cette seconde démonstration sont
reliés au théorème par des traits, et leur ordre de lecture est incertain ; on
les enchaîne selon l'argument. Deux mots de l'énoncé sont illisibles, et le
début en est biffé puis récrit (« image d'un morphisme » est ajouté).}

\textbf{Corollaire.} Si $f$ est fidèlement plat et quasi-compact, la topologie
de $S$ est la topologie quotient de celle de $S'$ : une partie $T$ de $S$ est
fermée si et seulement si $f^{-1}(T)$ l'est. En effet, si $f^{-1}(T)$ est
fermé, $T = f(f^{-1}(T))$ est l'image d'un morphisme quasi-compact, donc
$f^{-1}(\overline{T}) = \overline{f^{-1}(T)} = f^{-1}(T)$, et $T = \overline{T}$
par surjectivité.\footnote{« et quasi-compact » est ajouté au-dessus de la
ligne. La page écrit ensuite l'argument pour une partie fermée saturée de $S'$,
avec une ligne entière biffée.}

\textbf{Corollaire.} Si $f$ est fpqc, les suites
\[
\begin{array}{l}
F(S) \to F(S') \rightrightarrows F(S''), \\
O(S) \to O(S') \rightrightarrows O(S''), \\
LF_{\mathrm{rc}}(S) \to LF_{\mathrm{rc}}(S') \rightrightarrows LF_{\mathrm{rc}}(S'')
\end{array}
\]
sont exactes, où $F$, $O$, $LF_{\mathrm{rc}}$ désignent respectivement les
parties fermées, ouvertes, localement fermées rétrocompactes. Pour l'ensemble
$P$ de toutes les parties, l'exactitude vaut dès que $f$ est surjectif, parce
que $S'' \to S' \times_S S'$ (produit des ensembles sous-jacents) est
surjectif : une partie de $S'$ qui a même image réciproque par les deux
projections est saturée. Les deux premières suites s'en déduisent par le
corollaire précédent. Pour la troisième, on se ramène aux deux premières en
passant à l'adhérence et en utilisant le théorème : une partie $T$ qui est
l'image d'un morphisme quasi-compact est localement fermée si et seulement si
$f^{-1}(T)$ l'est.\footnote{La page porte, entre crochets et suivie d'un point
d'interrogation, la même suite pour les parties localement fermées sans
condition de rétrocompacité. Elle ne la tranche pas, et on ne la tranche pas ici :
c'est l'égalité $f^{-1}(\overline{T}) = \overline{f^{-1}(T)}$ qui manque pour
une partie quelconque. La même question revient à la page 5, pour les
sous-schémas. Le premier mot du dernier corollaire, lu « fermée » avec doute,
est probablement « ouverte », puisque « fermée » suit entre parenthèses.}

Ce qui sert dans toute la suite n'est pas la platitude en tant que telle, mais
le fait que $f$ et tous ses changements de base sont des applications
quotients. On dit aujourd'hui que $f$ est \emph{universellement submersif}. La
page note deux fois que cela vaut aussi pour $f$ surjectif, quasi-compact et
universellement générisant, ce que la démonstration du théorème
confirme.\footnote{Aux pages 3 et 4, entre parenthèses : « ou univ.
générisant », le second mot lu avec doute. Le nom « submersif » n'est pas sur
les pages.}

\pagerange{4}{4}
\subsection*{(1) et (2) : descente de propriétés de morphismes (page 4)}

Soit $g : X \to Y$ un $S$-morphisme et $g' : X' \to Y'$ son changement de base.
La question est : si $g'$ a une propriété, $g$ l'a-t-il ?\footnote{Sur toute la
page, le morphisme que l'on descend est appelé tantôt $g$, tantôt $f$, $f'$ ---
la lettre du morphisme de descente. Le contexte ne laisse pas de doute ; on écrit
$g$, $g'$ partout. La transcription le signale.}

\begin{itemize}
  \item Si $f$ est surjectif, $g$ est surjectif (resp. radiciel) si et seulement
  si $g'$ l'est.
  \item Si $f$ est surjectif et quasi-compact, $g$ est quasi-compact si et
  seulement si $g'$ l'est. Si $f$ est fpqc, il en va de même pour
  « de type fini », « de présentation finie », « localement de type fini »,
  « localement de présentation finie ».\footnote{La page écrit les deux cas en
  une phrase avec « resp. » ; on répartit les hypothèses comme elles doivent
  l'être : la quasi-compacité descend le long d'un morphisme surjectif
  quasi-compact, les conditions de finitude demandent la platitude fidèle.}
  \item \textbf{Lemme.} Si $A \to A'$ est fidèlement plat et $B$ une
  $A$-algèbre, $B$ est de type fini (resp. de présentation finie) si et
  seulement si $B \otimes_A A'$ l'est sur $A'$. Les mêmes arguments pour les
  modules donnent : si $f$ est fpqc et $F$ quasi-cohérent sur $S$, $F$ est de
  type fini (resp. de présentation finie, localement libre de type fini,
  localement libre de rang $n$) si et seulement si $F'$ l'est.
  \item Si $f$ est fidèlement plat, $g$ est plat (resp. fidèlement plat) si et
  seulement si $g'$ l'est.
\end{itemize}

Sous l'hypothèse que $f$ est fpqc, $Y' \to Y$ est une application quotient
universelle, ce qui donne la descente des propriétés topologiques (rubrique
(2)) : si $g'$ est ouvert (resp. fermé, générisant, un homéomorphisme sur son
image, un homéomorphisme), $g$ l'est ; donc $g$ est universellement ouvert
(resp. universellement fermé, un homéomorphisme universel) si et seulement si
$g'$ l'est.\footnote{La page annonce un « Th. » dont l'énoncé est laissé en
points de suspension, puis passe au corollaire. Une rubrique marginale
« Quelques contre-exemples » n'est suivie d'aucun.} Il suit que $g$ est séparé
(resp. propre) si et seulement si $g'$ l'est, et que $g$ est un isomorphisme si
et seulement si $g'$ l'est : $g$ est alors un homéomorphisme universel, et
$\mathcal{O}_Y \to g_*\mathcal{O}_X$ est un isomorphisme parce qu'il le devient
après le changement de base plat $Y' \to Y$.

Les immersions ouvertes, les immersions fermées et plus généralement les
immersions quasi-compactes descendent de même.\footnote{La page écrit
« $g$ une imm. (ouverte) ssi $g'$ l'est », le mot « ouverte » lu avec doute.
Pour une immersion quelconque, la descente se heurte à la question que la page
pose elle-même pour les parties localement fermées non rétrocompactes (page 3)
et pour les sous-schémas non rétrocompacts (page 5) ; on s'en tient aux cas
qu'elle établit.}

La page se termine par une liste, reliée par un trait aux corollaires : c'est
l'inventaire des propriétés qui descendent le long d'un morphisme fpqc.

\begin{itemize}
  \item surjectif, radiciel ;
  \item plat, fidèlement plat ;
  \item quasi-compact, quasi-séparé, qcqs, de type fini, de présentation
  finie, localement de type fini, localement de présentation finie ;
  \item séparé, propre, universellement fermé, universellement ouvert,
  homéomorphisme universel ;
  \item immersion fermée, immersion ouverte ;
  \item affine, entier, fini ;
  \item quasi-affine, quasi-fini ;
  \item lisse, non ramifié, étale ; formellement lisse, formellement
  étale.\footnote{La dernière ligne est ouverte par un crochet que rien ne
  ferme, et « formellement lisse » est suivi d'un point d'interrogation : la
  page ne se prononce pas, et nous non plus. Pour lisse, non ramifié et étale,
  la descente est établie dans EGA IV (§ 17). Le mot « net » de la page est le
  nom d'EGA pour « non ramifié ». Deux lignes de la liste sont biffées et
  illisibles, et la page porte « TSVP » en bas à droite.}
\end{itemize}

\pagerange{5}{6}
\subsection*{(3) Descente effective (pages 5 et 6)}

\textbf{Théorème.} Si $f : S' \to S$ est fpqc, $f$ est un morphisme de
descente effective, universellement, pour la catégorie fibrée des modules
quasi-cohérents.

Autrement dit --- et c'est ce que la page se donne à faire, « Expliquer ce que
cela signifie » ---, se donner un module quasi-cohérent $F$ sur $S$ revient au
même que se donner un module quasi-cohérent $F'$ sur $S'$ muni d'un isomorphisme
$\varphi : p_1^* F' \simeq p_2^* F'$ sur $S''$ qui vérifie la condition de
cocycle sur $S' \times_S S' \times_S S'$ ; et cela reste vrai après tout
changement de base $T \to S$.\footnote{La page n'écrit pas cette explication ; elle
en note seulement la nécessité. Celle qu'on donne est la définition usuelle.}

Les corollaires en découlent en appliquant le théorème à des sous-objets.

\begin{itemize}
  \item Pour $F$ quasi-cohérent sur $S$, la suite des ensembles de
  sous-modules quasi-cohérents $\mathrm{Ssf}(F) \to \mathrm{Ssf}(F')
  \rightrightarrows \mathrm{Ssf}(F'')$ est exacte.\footnote{La page emploie la même
  abréviation $\mathrm{Ssf}$ pour les sous-modules d'un module et pour les
  sous-schémas fermés d'un schéma. Un premier énoncé, cerclé et barré, le
  démontrait en se ramenant au cas affine ; son croquis en marge
  ($N' \cap M = N_0$) indique qu'il s'agit bien de sous-modules.}
  \item Les suites analogues pour les sous-schémas fermés ($\mathrm{Ssf}$),
  ouverts ($\mathrm{Sso}$) et rétrocompacts, c'est-à-dire dont l'immersion est
  quasi-compacte ($\mathrm{Ssrc}$), sont exactes. La page demande si l'on peut
  descendre les sous-schémas non rétrocompacts ; la question reste ouverte sur
  la page.
  \item $f$ est un morphisme de descente, universellement, pour la catégorie
  fibrée des schémas : la suite
  \[
  \operatorname{Hom}_S(X, Y) \to \operatorname{Hom}_{S'}(X', Y')
  \rightrightarrows \operatorname{Hom}_{S''}(X'', Y'')
  \]
  est exacte. Lorsque $Y$ est quasi-séparé sur $S$, cela suit du corollaire
  précédent appliqué aux graphes, qui sont des sous-schémas rétrocompacts de
  $X \times_S Y$ ; en général on se ramène à $S$ et $Y$ affines. En langage
  d'aujourd'hui : les foncteurs représentables sont des faisceaux pour la
  topologie fpqc.
  \item Une immersion fermée (resp. ouverte) descend ; les algèbres
  quasi-cohérentes descendent.
\end{itemize}

\textbf{Théorème.} Si $f$ est fpqc, $f$ est un morphisme de descente
effective, universellement, pour la catégorie fibrée des schémas
\emph{affines} sur la base.

\textbf{Corollaire.} $g$ est affine (resp. quasi-affine, entier, fini, fini
localement libre) si et seulement si $g'$ l'est. Pour « quasi-affine », on
utilise que $g$ est quasi-affine si et seulement si $g$ est qcqs et
$X \to \operatorname{Spec} g_*(\mathcal{O}_X)$ est une immersion
ouverte.\footnote{La page écrit « immersion ». Quand $g$ est qcqs, le morphisme
$X \to \operatorname{Spec} g_*(\mathcal{O}_X)$ est quasi-compact, et c'est une
immersion si et seulement si c'est une immersion ouverte : les deux énoncés
sont équivalents.} Pour « fini » et « entier », on est ramené au lemme de la
page 6 : si $A \to A'$ est fidèlement plat, une $A$-algèbre $B$ est finie (resp.
entière) si et seulement si $B \otimes_A A'$ l'est sur $A'$. « Fini est une
histoire de module » ; entier s'y ramène en regardant les sous-algèbres
$A[x]$, $x \in B$.

\pagerange{6}{6}
\subsection*{(4) Descente de l'amplitude, et (5) la rubrique vide (page 6)}

Soient $g : X \to Y$ un $S$-morphisme qcqs, $L$ un faisceau inversible sur $X$,
et $f : S' \to S$ fpqc. Alors $L$ est très ample (resp. ample) relativement à
$Y$ si et seulement si $L'$ l'est relativement à $Y'$.

La démonstration de la page est la suivante. Comme $g$ est qcqs, $g_*(L)$ est
quasi-cohérent et $g_*(L)' \simeq g'_*(L')$. « Très ample » signifie que
$g^* g_*(L) \to L$ est surjectif, ce qui définit
$X \to \mathbf{P}(g_*(L))$, et que ce morphisme est une immersion ; la
surjectivité descend, et l'immersion, qui est quasi-compacte, aussi. Si $g$ est
de type fini, « ample » se ramène à « très ample », localement sur $Y$. En
général, on utilise le critère : pour $Y$ quasi-compact, $L$ est ample
relativement à $Y$ si et seulement si, pour tout module quasi-cohérent de type
fini $F$ sur $X$, $g^* g_*(F \otimes L^{\otimes n}) \to F \otimes L^{\otimes n}$
est surjectif pour $n$ assez grand.\footnote{La page écrit
$f^*(f_*(F(n)))$ et « $F$ … sur $S$ » : ce sont les lettres du morphisme de
descente et de la base, mis pour le morphisme structural $g$ et pour $X$. On
corrige.} On retrouve la descente des morphismes quasi-affines, puisque $g$ est
quasi-affine si et seulement si $\mathcal{O}_X$ est ample relativement à $Y$.

\textbf{NB.} Parmi les propriétés usuelles de morphismes, les seules de la liste
qui ne descendent pas sont « projectif » et « quasi-projectif ».\footnote{La
page dit « les seuls exemples », avec un mot illisible ; on restreint
l'affirmation à ce qu'elle peut vouloir dire. La raison est dans ce qui précède :
l'amplitude d'un faisceau \emph{donné} descend, mais un faisceau inversible sur
$X'$ ne provient pas en général d'un faisceau sur $X$. La projectivité n'est
d'ailleurs pas même locale sur la base pour la topologie de Zariski, comme le
montre un exemple classique de Hironaka.}

La page se clôt sur un titre : « (5) Questions d'effectivité de données de
descente ». Rien n'est écrit dessous, et le bas de la feuille est
vide.\footnote{On ne sait pas ce qu'il comptait y mettre. Les réponses
connues aujourd'hui sont celles de l'exposé VIII de SGA 1 : une donnée de
descente est effective pour les schémas affines sur la base, et plus
généralement pour les schémas quasi-projectifs munis d'un faisceau ample
compatible avec la donnée ; elle ne l'est pas en général, et c'est l'une des
origines des espaces algébriques.}

\pagerange{7}{9}
\section*{II. Trois classes de morphismes (pages 7 à 9)}

\pagerange{7}{7}
\subsection*{Morphismes radiciels, homéomorphismes universels, morphismes universellement ouverts (page 7)}

Soit $f : X \to Y$ un morphisme de schémas. Les conditions suivantes sont
équivalentes :\footnote{La page numérote ces conditions a) à d). On les
renumérote pour réserver a), b), c) aux règles de stabilité de la page 10,
auxquelles la même page renvoie deux lignes plus bas.}

\begin{itemize}
  \item[(i)] $f$ est universellement injectif ;
  \item[(ii)] pour tout $y \in Y$, la fibre $X_y$ a au plus un point $x$, et
  $k(x)$ est une extension radicielle de $k(y)$ ;
  \item[(iii)] pour tout corps $K$, l'application $X(K) \to Y(K)$ est
  injective ;
  \item[(iv)] le morphisme diagonal $\delta_f : X \to X \times_Y X$ est
  surjectif.
\end{itemize}

Ces morphismes sont dits \emph{radiciels}. Une immersion est radicielle ; donc
le morphisme diagonal, qui est une immersion, l'est. La classe des morphismes
radiciels satisfait les trois règles a), b), c) : stabilité par changement de
base, par composition, et si $gf$ est radiciel, $f$ l'est.

\textbf{Proposition.} Un morphisme entier, radiciel et surjectif est un
homéomorphisme universel. La page ajoute qu'on peut démontrer la réciproque,
« mais c'est nettement plus difficile », et ne la démontre pas.\footnote{La
réciproque est vraie : un homéomorphisme universel est entier, radiciel et
surjectif (EGA IV, § 18.12). « surjectif » est ajouté au-dessus de la ligne dans
l'énoncé.}

Exemples : si $k'/k$ est une extension radicielle (par exemple la clôture
parfaite), $X_{k'} \to X$ est entier, radiciel et surjectif, donc un
homéomorphisme universel ; de même la normalisation d'une courbe à point de
rebroussement, $t \mapsto (t^2, t^3)$ de la droite vers la cubique
$y^2 = x^3$.

\emph{Morphismes universellement ouverts.} Un morphisme ouvert est
générisant ; la réciproque vaut pour un morphisme localement de présentation
finie. Donc un morphisme universellement ouvert est universellement
générisant, et réciproquement s'il est localement de présentation finie. Les
deux classes sont stables par changement de base et par
composition.\footnote{La page écrit « satisfont 1), 2) », qui renvoie à ces
deux règles. Le mot « générisant » y est chaque fois lu avec doute ; c'est le
seul qui donne un sens mathématique juste aux phrases.}

\pagerange{8}{8}
\subsection*{Morphismes quasi-affines (page 8)}

Un schéma est \emph{quasi-affine} s'il est quasi-compact et isomorphe à un
ouvert d'un schéma affine. Pour un schéma $X$ qcqs, d'anneau
$A = \Gamma(X, \mathcal{O}_X)$, les conditions suivantes sont équivalentes :

\begin{itemize}
  \item[(a)] $X$ est quasi-affine ;
  \item[(b)] $X \to \operatorname{Spec} A$ est une immersion ouverte ;
  \item[(b')] $X \to \operatorname{Spec} A$ est un homéomorphisme sur son
  image ;
  \item[(c)] $\mathcal{O}_X$ est ample, c'est-à-dire : les $X_f$
  ($f \in A$) forment une base de la topologie ; ou : les $X_f$ affines en
  forment une ; ou : tout module quasi-cohérent est engendré par ses sections
  globales ; ou : tout idéal quasi-cohérent l'est ;
  \item[(d)] il existe un faisceau inversible $L$ tel que $L$ et $L^{-1}$
  soient amples ;
  \item[(d')] tout faisceau inversible est ample.\footnote{En (d), « inversible »
  est écrit au-dessus d'un « ample » biffé.}
\end{itemize}

Les morphismes quasi-affines sont stables par changement de base et par
composition ; une immersion quasi-compacte est quasi-affine, donc la diagonale
d'un morphisme quasi-séparé l'est ; si $gf$ est quasi-affine et $g$
quasi-séparé, $f$ est quasi-affine.

\textbf{Théorème.} Un morphisme séparé et quasi-fini (de type fini, à fibres
finies) est quasi-affine, donc quasi-projectif. C'est une forme du
\emph{Main Theorem} de Zariski.\footnote{Le nom n'est pas sur la page.}

Les définitions sont celles d'EGA II :\footnote{Elles diffèrent de celles de
plusieurs manuels d'aujourd'hui, où « projectif » veut dire « immersion fermée
dans un $\mathbf{P}^n_Y$ ». Dans la convention d'EGA, $\mathbf{P}(E)$ est pris
pour un module quasi-cohérent de type fini $E$ quelconque.} $f$ est
\emph{quasi-projectif} s'il est de type fini et s'il existe un faisceau
inversible $f$-ample ; \emph{propre} s'il est séparé, de type fini et
universellement fermé ; et, lorsque $Y$ est quasi-compact, $f$ est
\emph{projectif} si et seulement s'il est propre et quasi-projectif, ou encore
universellement fermé et quasi-projectif.\footnote{La page écrit
« projectif $=$ propre $+$ quasi-projectif » sans condition sur $Y$. L'égalité
demande que $Y$ soit quasi-compact : sans cela un faisceau ample relatif peut
exiger une infinité de sections, et le module $E$ de la définition n'est plus
de type fini.}

Le tableau encadré de la page se lit ainsi, chaque flèche étant une
implication :

\begin{tikzcd}
\text{imm. fermée} \arrow[r, Rightarrow] \arrow[d, Rightarrow] & \text{fini} \arrow[r, Rightarrow] \arrow[d, Rightarrow] & \text{entier} \arrow[r, Rightarrow] \arrow[d, Rightarrow] & \text{affine} \arrow[d, Rightarrow] \\
\text{imm. quasi-compacte} \arrow[d, Rightarrow] & \text{projectif} \arrow[r, Rightarrow] \arrow[d, Rightarrow] & \text{univ. fermé} & \text{séparé} \\
\text{quasi-affine de t.f.} \arrow[r, Rightarrow] & \text{quasi-projectif} &  &
\end{tikzcd}

avec en outre : projectif $\Rightarrow$ propre $\Rightarrow$ universellement
fermé ; propre $\Rightarrow$ séparé ; quasi-projectif $\Rightarrow$ séparé ;
fini $\Rightarrow$ affine de type fini $\Rightarrow$ quasi-affine de type
fini.\footnote{Sur la page, « propre » est intercalé sur la flèche qui va de
« proj » à « univ. fermé » ; on l'a sorti du dessin. Une flèche barrée, venant
de la marge, entre dans le cadre ; « affine de t.f. » est écrit contre une
flèche oblique. Toutes les implications sont justes avec les définitions
d'EGA ; celle de « fini » vers « projectif » est la moins évidente.}

\textbf{Théorème.} Un morphisme est fini si et seulement s'il est propre et à
fibres finies, si et seulement s'il est projectif et à fibres
finies.\footnote{La page qualifie l'énoncé de « cas noethérien », lu avec
doute. L'énoncé vaut sans hypothèse noethérienne (EGA IV, § 18.12) ; le cas
noethérien est dans EGA III. Une seconde assertion (b), sur les morphismes
séparés à fibres finies, est entièrement biffée.}

\pagerange{9}{9}
\subsection*{Faisceaux amples (page 9)}

La marge porte « EGA II 4.5 », et la page résume en effet ce paragraphe
d'EGA. Soient $X$ un schéma qcqs et $L$ un faisceau inversible ; pour
$f \in \Gamma(X, L^{\otimes n})$, $n \geqslant 1$, on note $X_f$ l'ouvert où
$f$ ne s'annule pas. Les conditions suivantes sont équivalentes :

\begin{itemize}
  \item[(a)] les $X_f$ forment une base de la topologie de $X$ ;
  \item[(a')] les $X_f$ affines en forment une ;
  \item[(b)] pour tout module quasi-cohérent de type fini $F$,
  $F \otimes L^{\otimes n}$ est engendré par ses sections globales pour $n$
  grand ; autrement dit $F$ est quotient d'une somme de copies de
  $L^{\otimes -n}$ ;
  \item[(b')] la même chose pour les idéaux quasi-cohérents de type fini.
\end{itemize}

Si $X$ est au-dessus d'un schéma affine $S$, il suffit pour cela qu'une
puissance $L^{\otimes n}$, $n > 0$, soit très ample relativement à $S$ ; la
condition est nécessaire si $X$ est de type fini sur $S$, et alors
$L^{\otimes n}$ est très ample pour tout $n$ assez grand. On dit que $L$ est
\emph{ample} s'il satisfait ces conditions ; si $f : X \to Y$ est qcqs, $L$ est
\emph{$f$-ample} s'il existe un recouvrement de $Y$ par des ouverts affines
$Y_i$ au-dessus desquels $L$ est ample ; c'est alors vrai au-dessus de tout
ouvert affine. Si une puissance de $L$ est très ample relativement à $Y$, $L$
est $f$-ample ; la réciproque vaut si $f$ est de type fini et $Y$
quasi-compact.

\textbf{Proposition.} Soient $f : X \to Y$ et $g : Y \to Z$ qcqs, avec $Z$
quasi-compact, $L$ ample relativement à $Y$ et $M$ ample relativement à $Z$.
Il existe $n_0$ tel que $L \otimes f^*(M^{\otimes n})$ soit ample
relativement à $Z$ pour tout $n \geqslant n_0$. De même pour « très
ample ».\footnote{La page suppose seulement $g$ quasi-compact. Pour que le
même $n_0$ convienne partout, il faut que $Z$ soit quasi-compact ; sans cela
on n'a l'énoncé qu'au-dessus de chaque ouvert quasi-compact de $Z$. « idem avec
très amples » est sur la page ; dans ce cas on peut prendre $n = 1$.}

\pagerange{10}{12}
\section*{III. Plans de cours (pages 10 à 12)}

Trois pages de plans, les deux premières à l'encre, la troisième au crayon,
surchargées de flèches, d'accolades et de reprises. On n'en donne pas une
lecture ligne à ligne --- la transcription le fait ---, mais ce qu'elles fixent :
une manière d'organiser les classes de morphismes, et deux ordres de
traitement.

\pagerange{10}{11}
\subsection*{Les règles de stabilité a), b), c) (pages 10 et 11)}

Soit $\mathcal{M}$ une classe de morphismes de schémas contenant les
isomorphismes.\footnote{La page écrit, après la règle a), « ($\Rightarrow$ tout
isom. $\in \mathcal{M}$) ». La stabilité par changement de base ne donne cela
que si $\mathcal{M}$ contient déjà des identités ; on le suppose, comme le fait
toute la suite de la page.}

\begin{itemize}
  \item[a)] Si $f \in \mathcal{M}$, tout changement de base $f'$ de $f$ est
  dans $\mathcal{M}$.
  \item[b)] Si $f, g \in \mathcal{M}$, alors $gf \in \mathcal{M}$.
  \item[c)] Si $gf \in \mathcal{M}$, alors $f \in \mathcal{M}$.
\end{itemize}

De a) et b) suit que $f_1 \times_S f_2 \in \mathcal{M}$ si
$f_1, f_2 \in \mathcal{M}$, et la règle
\[
(*) \qquad gf \in \mathcal{M} \text{ et } \Delta_g \in \mathcal{M}
\ \Longrightarrow\ f \in \mathcal{M} .
\]
En effet, pour $f : X \to Y$ et $g : Y \to S$, $f$ se factorise en
$X \xrightarrow{\Gamma_f} X \times_S Y \xrightarrow{\mathrm{pr}_2} Y$ ; le
graphe $\Gamma_f$ est le changement de base de $\Delta_g : Y \to Y \times_S Y$
par $f \times \mathrm{id}_Y$, et $\mathrm{pr}_2$ celui de $gf : X \to S$ par
$g$.

\begin{tikzcd}
X \arrow[r, "\Gamma_f"] \arrow[d, "f"'] & X \times_S Y \arrow[d, "f \times \mathrm{id}_Y"] \\
Y \arrow[r, "\Delta_g"'] & Y \times_S Y
\end{tikzcd}

Sous a) et b), pour un morphisme $g : Y \to S$, les conditions suivantes sont
équivalentes :

\begin{itemize}
  \item[(i)] $\Delta_g \in \mathcal{M}$ ;
  \item[(ii)] pour tout $S' \to S$, toute section de $Y' \to S'$ est dans
  $\mathcal{M}$ ;
  \item[(iii)] pour tout couple de $S$-morphismes $u, v : X \rightrightarrows Y$,
  le noyau $\operatorname{Ker}(u, v) \to X$ est dans $\mathcal{M}$ ;
  \item[(iv)] pour tout $S$-morphisme $f : X \to Y$, le graphe $\Gamma_f$ est
  dans $\mathcal{M}$ ;
  \item[(v)] pour tout $S' \to S$ et tout $S'$-morphisme $f : X \to Y'$,
  $g'f \in \mathcal{M}$ entraîne $f \in \mathcal{M}$.
\end{itemize}

Et c) équivaut, sous a) et b), à : toutes les diagonales sont dans
$\mathcal{M}$.\footnote{La page écrit la dernière équivalence
« $gf \in \mathcal{M} \Rightarrow f \in \mathcal{M}$ » pour $g$ seul. Pour
qu'elle entraîne (i), il faut la demander aussi pour les changements de base de
$g$ : c'est ce qui permet d'y prendre pour $f$ une section, comme en (ii). Une
variante entre crochets, pour les noyaux, suppose de plus la propriété
« locale en bas et vraie au cas affine ».}

Les plans passent alors les classes en revue, en notant pour chacune les règles
qu'elle satisfait. Morphismes quasi-compacts, de type fini : a) et b).
Localement de type fini, immersions, morphismes séparés : a), b), c) ; si $f$
est localement de type fini, sa diagonale est localement de présentation finie.
Morphismes \emph{séparés} : la diagonale est fermée, ou les noyaux de couples
sont fermés, ou encore vaut le principe de prolongement des identités --- deux
$S$-morphismes d'un schéma réduit $Z$ vers $X$ qui coïncident sur un ouvert
dense sont égaux.\footnote{La page 11 écrit « pour $X$ réduit » ; c'est la
source $Z$ qui doit l'être, comme le dit la page 10 (« morphismes des schémas
réduits dedans »).} Morphismes \emph{quasi-séparés} : la diagonale est
quasi-compacte ; sur une base quasi-séparée, cela revient à dire que
l'intersection de deux ouverts quasi-compacts est quasi-compacte. Morphismes
\emph{qcqs} : a), b), et si $gf$ est qcqs avec $g$ quasi-séparé, $f$ l'est ;
si $f$ est qcqs et $F$ quasi-cohérent, $f_*(F)$ et les $R^i f_*(F)$ sont
quasi-cohérents.\footnote{La page 10 écrit « Si $f$ qu.-coh. » pour « si $f$
cohérent » ; la transcription le signale.} Morphismes de présentation finie
(localement de présentation finie et qcqs) : a), b), et si $gf$ est de
présentation finie avec $g$ de type fini et quasi-séparé, $f$ l'est. Toutes ces
notions sont locales sur la base.

Au milieu de la page 10, un NB mérite d'être rapporté pour lui-même : tous les
schémas envisagés par « les Anciens » étaient de type fini et séparés sur des
corps ; ceux qu'on envisage couramment pourraient être supposés séparés,
pratiquement tous ceux qu'on rencontre sont quasi-séparés, et beaucoup sont
quasi-compacts --- « mais il y en a d'importants qui ne le sont pas ».

\pagerange{10}{11}
\subsection*{Deux programmes (pages 10 et 11)}

Le bas de la page 10 donne un plan en quatre temps :

\begin{itemize}
  \item[a)] Traversée du désert ;
  \item[b)] Passage au quotient, et langage faisceautique (descente
  fidèlement plate) ;
  \item[c)] Notion « géométrique » sur un corps ;
  \item[d)] Ensembles constructibles, passage à la limite.
\end{itemize}

C'est à peu près, sans que la page le dise, l'ordre de la première moitié
d'EGA IV.\footnote{Rapprochement de nous : EGA IV traite d'abord des conditions
de finitude relatives, puis de la platitude et de la descente, puis des
propriétés « géométriques » par changement de corps, puis des limites
projectives et des propriétés constructibles.}

La page 11 dresse un autre programme, en trois blocs. I : morphismes
quasi-compacts, de type fini, localement de présentation finie ; immersions
ouvertes et fermées ; morphismes séparés ; les règles de stabilité. II :
schémas noethériens, localement noethériens, artiniens ; irréductibles,
connexes ; intègres ; normaux, réguliers ; puis les schémas locaux --- un
schéma quasi-compact qui n'a qu'un point fermé est le spectre d'un anneau
local, puisque tout ouvert affine contenant ce point contient tous les
autres.\footnote{La page définit le schéma local par « $\exists !$ pt fermé, et
quasi-compact » ; la justification entre tirets est de nous.} III : les
morphismes affines, quasi-séparés, qcqs, de présentation finie ; finis,
entiers, quasi-finis ; projectifs, quasi-projectifs, quasi-affines ; radiciels,
homéomorphismes universels ; propres. Des lettres cerclées (a) à (d) y fixent un
ordre de traitement qui n'est pas celui de la liste : projectifs, puis propres,
puis finis, puis radiciels. Une rubrique « Oubli » rappelle trois points : la
cohérence du faisceau structural d'un schéma localement noethérien, la
structure réduite induite sur un fermé, et qu'un schéma localement noethérien
est quasi-séparé. Une bulle ajoute : « on verra tôt les schémas non
séparés ». Suivent le prolongement d'un module quasi-cohérent, ou d'un
sous-module, depuis un ouvert rétrocompact, et une accolade où « plat,
fidèlement plat » est encadré.\footnote{La page est traversée de longues barres
obliques, comme pour marquer le bloc comme fait ; elle reste lisible. Une
seconde bulle ne se lit qu'en partie.}

\pagerange{12}{12}
\subsection*{Le plan au crayon : schémas comme foncteurs, et $P(M)$ (page 12)}

La page 12 est la suite d'un plan à rubriques cerclées, dont se lisent (d) à
(h). On y voit apparaître le point de vue que les pages 14 et 16 développeront.

(d) Schémas de type fini sur $S$.\footnote{La page les définit par
« quasi-comp. et qu.-sép. sur $S$ et loc. de type fini sur $S$ ». La
convention d'EGA ne demande pas « quasi-séparé » ; la page adopte une variante,
qu'on signale sans la corriger, puisque c'est une définition.} Puis : « Schémas
comme foncteurs : avantages respectifs des arguments anneaux (schémas affines)
et des arguments schémas (vrais espaces localement annelés) ». La rubrique (e),
biffée, annonçait un « corollaire du théorème de représentabilité ».

(f) Le schéma $P(M)$, pour $M$ quasi-cohérent sur $S$. Il représente un
foncteur « remarquable » sur la catégorie des espaces localement annelés
au-dessus de $S$ : celui qui associe à $T$ l'ensemble des quotients inversibles
de $M_T$. La page ajoute entre parenthèses « univ. des topos loc. annelés sur
$S$ ».\footnote{Ce point de vue sera développé par Monique Hakim, élève de
Grothendieck, dans \emph{Topos annelés et schémas relatifs} (1972). La page ne
cite personne ; le rapprochement est de nous.} Ses dépendances : par
changement de base, fonctorielle en $S$, et en $M$ pour les épimorphismes. Les
schémas quasi-projectifs sur $S$ affine se plongent dans un $P(M)$, avec $M$ de
type fini.

Relation avec le fibré vectoriel $V(M) = \operatorname{Spec}(\operatorname{Sym} M)$,
dont les points à valeurs dans $T$ sont les formes linéaires
$M_T \to \mathcal{O}_T$. Soit $V(M)^{*} \subset V(M)$ le sous-foncteur ouvert
des formes \emph{surjectives}. Le morphisme $V(M)^{*} \to P(M)$ qui envoie une
forme surjective sur le quotient qu'elle définit est un épimorphisme de
faisceaux pour la topologie de Zariski --- tout quotient inversible admet
localement une telle forme, là où il est trivial. Le groupe
$\mathbb{G}_{m,S}$ opère sur $V(M)$ par homothéties, laisse $V(M)^{*}$ stable
et y opère librement, et
\[
P(M) \simeq V(M)^{*} / \mathbb{G}_m \qquad \text{(isomorphisme de faisceaux de Zariski).}
\]
C'est la présentation de l'espace projectif comme quotient de « l'espace
affine privé de l'origine » par les homothéties, écrite pour un module
quelconque.\footnote{La page écrit que $\mathbb{G}_m$ « opère sur $V(M)^{*}$,
laissant $V(M)^{*}$ stable » : la première occurrence est à lire $V(M)$. Le
quotient est un quotient de faisceaux, non de schémas : c'est ce que dit la
parenthèse de la page.}

(g) Limites projectives finies dans la catégorie des schémas. (h) Spectre d'une
algèbre quasi-cohérente, et morphismes affines.

\pagerange{14}{14}
\section*{IV. Les schémas comme foncteurs (pages 14 et 16)}

\subsection*{Morphismes vers un schéma affine, et le diagramme des faisceaux (page 14)}

\textbf{Proposition.} Soient $\mathcal{X}$ un espace localement annelé au-dessus
de $k$ et $A$ une $k$-algèbre. L'application
\[
\operatorname{Hom}_k(\mathcal{X}, \operatorname{Spec} A) \longrightarrow
\operatorname{Hom}_{k\text{-alg}}\bigl(A, \Gamma(\mathcal{X}, \mathcal{O}_{\mathcal{X}})\bigr)
\]
est bijective.

La page construit l'inverse. Soit $u : A \to \Gamma(\mathcal{X}, \mathcal{O}_{\mathcal{X}})$.
Pour $x \in \mathcal{X}$, on pose $\varphi(x) = $ le noyau du composé
$A \to \Gamma(\mathcal{X}, \mathcal{O}_{\mathcal{X}}) \to k(x)$, qui est un idéal
premier. Pour $f \in A$, $\varphi^{-1}(Y_f)$ est l'ensemble des $x$ où $u(f)_x$
est inversible, c'est-à-dire l'ouvert $\mathcal{X}_{u(f)}$ : $\varphi$ est
continue. Les homomorphismes
$A_f \to \Gamma(\mathcal{X}_{u(f)}, \mathcal{O}_{\mathcal{X}})$, qui rendent $f$
inversible, sont compatibles aux restrictions et définissent un morphisme
d'espaces annelés,\footnote{La page écrit
$\Gamma(U_f, \mathcal{O}_{\mathcal{X}})$ ; l'ouvert de $\mathcal{X}$ dont il
s'agit est $\varphi^{-1}(U_f) = \mathcal{X}_{u(f)}$.} qui est local sur les
fibres parce que $\mathcal{X}$ est localement annelé.\footnote{Cette dernière
phrase est sur la page en partie illisible ; on donne la raison usuelle.}

\emph{Conséquence.} La catégorie des $k$-schémas est la catégorie des
schémas au-dessus de $S = \operatorname{Spec} k$ ; mais la catégorie
$\mathrm{Sch}_{/S}$ est importante sans que $S$ soit nécessairement affine.

Le grand diagramme de la page met en place trois catégories et leurs
catégories de préfaisceaux et de faisceaux de Zariski :
\[
\mathrm{Aff}_k \simeq (\mathrm{Alg}_k)^{\circ} \subset \mathrm{Sch}_k \subset \mathrm{Locan}_k ,
\]
où $\mathrm{Locan}_k$ est la catégorie des espaces localement annelés au-dessus
de $k$, chaque inclusion étant pleine. Les préfaisceaux sur $\mathrm{Aff}_k$
sont les foncteurs covariants $\mathrm{Alg}_k \to \mathrm{Ens}$. La restriction
des préfaisceaux de $\mathrm{Sch}_k$ à $\mathrm{Aff}_k$ n'est pas pleinement
fidèle, et la page le note.\footnote{La page écrit « (pas plein fid.) » sous la
flèche de restriction des préfaisceaux ; « fid. » y est lu avec doute. L'énoncé
qu'on donne est juste quelle que soit cette lecture.} En revanche la restriction des \emph{faisceaux} de
Zariski est une équivalence :
\[
\widetilde{\mathrm{Sch}_k} \xrightarrow{\ \sim\ } \widetilde{\mathrm{Aff}_k}
\simeq \{\text{faisceaux de Zariski } \mathrm{Alg}_k \to \mathrm{Ens}\} ,
\]
parce que tout schéma est recouvert par des schémas affines.\footnote{La page
trace la flèche de restriction sans l'étiqueter « $\simeq$ ». L'équivalence est
le lemme de comparaison des sites (SGA 4, exposé III), aux précautions d'univers
près ; c'est elle qui permet, page 16, de parler indifféremment de foncteurs sur
les schémas et de foncteurs sur les algèbres.}

La page pose alors la question : comment exprimer en algèbre commutative
qu'un foncteur $F : \mathrm{Alg}_k \to \mathrm{Ens}$ est un faisceau de
Zariski ? La réponse est : pour toute $k$-algèbre $A$ et toute famille finie
$(f_i)$ d'éléments engendrant l'idéal unité, la suite
\[
F(A) \to \prod_i F(A_{f_i}) \rightrightarrows \prod_{i,j} F(A_{f_i f_j})
\]
est exacte.\footnote{La page pose la question sans y répondre ; la réponse est
de nous. Elle contient le cas des sommes finies, puisque
$\operatorname{Spec}(A \times B)$ est recouvert par les deux ouverts
$\operatorname{Spec} A$ et $\operatorname{Spec} B$.} Un ajout au crayon rappelle
que $\widehat{C}_{/S} \simeq \widehat{C_{/S}}$ pour $C = \mathrm{Aff}_k$ :
préfaisceaux au-dessus d'un objet et préfaisceaux sur la catégorie
au-dessus.\footnote{La transcription écrit les deux membres de la même façon,
$\widehat{C}_{/S} \approx \widehat{C}_{/S}$ ; c'est l'équivalence classique
qu'on donne, la seule qui ait un sens ici. Le signe $\approx$ est souligné d'un
trait épais.} Suivent deux notes : les sommes s'interprètent en termes de
faisceaux (sup de sous-faisceaux), et les limites projectives finies ne posent
pas de difficulté.

\pagerange{16}{16}
\subsection*{Le critère de représentabilité (page 16)}

Un sous-foncteur $F_i \subset F$ est \emph{ouvert} si, pour tout schéma $T$ et
tout $T \to F$, le produit fibré $T \times_F F_i$ est représenté par un ouvert
de $T$. Une famille de sous-foncteurs ouverts \emph{recouvre} $F$ si, pour
tout $T \to F$, les ouverts correspondants recouvrent $T$ ; c'est ce que la page
écrit $F = \operatorname{Sup} F_i$, le sup étant pris dans la catégorie des
faisceaux.

\textbf{Proposition.} Un foncteur $F : (\mathrm{Sch}_{/S})^{\circ} \to \mathrm{Ens}$
est représentable si et seulement si (a) $F$ est un faisceau de Zariski, et
(b) $F$ est recouvert par une famille de sous-foncteurs ouverts représentables.

\textbf{Corollaire.} On peut exiger en (b) que les $F_i$ soient des schémas
affines. En particulier, pour $S = \operatorname{Spec} k$, un foncteur
$\mathrm{Alg}_k \to \mathrm{Ens}$ est un schéma si et seulement si c'est un
faisceau de Zariski qu'on peut recouvrir par des sous-foncteurs ouverts
affines.\footnote{La page écrit : « Pour qu'un foncteur $\mathrm{Alg}_k \to
(\mathrm{Ens})$ soit un faisceau il f. et il s. que ce soit un faisceau, et
qu'on le puisse recouvrir… » ; le premier « faisceau » est un lapsus pour
« schéma », que la transcription relève (\emph{sic}).}

Et la page conclut : « Ainsi, la notion de schéma est explicitée sans recours à
la notion d'espaces \emph{topologiques}, faisceau d'anneaux, espaces loc.
annelés ». C'est la définition des schémas par leur \emph{foncteur des points},
qu'on trouve ainsi rédigée dans le traité de Demazure et Gabriel (1970).\footnote{La
page ne cite rien ; le rapprochement est de nous et ne dit rien de l'ordre des
dates.}

Applications annoncées : (a) représentabilité du foncteur $P(E)$, $E$ localement
libre sur $S$ --- par exemple un espace vectoriel de dimension finie sur un
corps --- ; (b) représentabilité de $\operatorname{Spec}$ d'une algèbre ; (c)
produits fibrés, et plus généralement limites projectives finies ; (d) limites
projectives filtrantes de schémas affines au-dessus d'un schéma. Une bulle en
marge ajoute : sommes disjointes, et disjointes universelles.

\pagerange{22}{24}
\section*{V. Le cours : prérequis et plan (pages 22 à 24, 26 à 28)}

\pagerange{22}{22}
\subsection*{La page anglaise (page 22)}

La page s'ouvre en anglais : « Prerequisites for the courses I and II ». Le cours
I est une introduction à la géométrie algébrique, qui suppose la théorie des
catégories et de l'algèbre de base, y compris un peu d'algèbre commutative --- la
page renvoie à Bourbaki et promet de rappeler la plupart des faits utilisés.
Le cours II est une introduction aux groupes algébriques, qui suppose les
notions de base de la géométrie algébrique. Le sujet, dit la page, est
« messy » ; on ne pourra en donner, dans ce cours d'été, qu'un survol de
quelques points clés, en esquissant des démonstrations le plus souvent
directes, et un travail personnel sera nécessaire à qui n'a pas déjà une bonne
formation. Et, pour finir : « The content of I are prerequisites for II! ».

Le plan, en français, commence aussitôt dessous.

\subsection*{(1)--(3) Équations, et le point de vue géométrique (page 22)}

Soient $k$ un anneau, $I$ et $J$ des ensembles finis ou infinis, et
$P_I = k[(T_i)_{i \in I}]$.\footnote{La page écrit $P[(T_i)_{i \in I}]$, pour
$k[\ldots]$.} Un point $x \in k^I$ est la même chose qu'un homomorphisme de
$k$-algèbres $\varepsilon_x : P_I \to k$, et $f(x) = \varepsilon_x(f)$. On
considère un système d'équations $F_j(T) = 0$, $j \in J$.

Il y a deux points de vue. Le point de vue \emph{arithmétique} étudie les
solutions dans $k$ lui-même, et il dépend fortement de la nature de $k$. Le
point de vue \emph{géométrique}, issu de la géométrie sur $\mathbb{R}$ et sur
$\mathbb{C}$, étudie les solutions dans toutes les $k$-algèbres $k'$, la manière
dont elles varient, et les propriétés du système qui sont stables par
changement de $k'$. On introduit donc le foncteur
\[
\mathbb{E}^I_k : \mathrm{Alg}_k \to \mathrm{Ens}, \qquad k' \longmapsto k'^I ,
\]
et son sous-foncteur $V_S \subset \mathbb{E}^I_k$ des solutions du système
$S = (F_j)$. On peut étudier $V_S$ indépendamment de son plongement, ou étudier le
plongement lui-même.

\pagerange{23}{23}
\subsection*{(4)--(5) Les schémas affines comme foncteurs représentables (page 23)}

Un sous-foncteur fermé de $\mathbb{E}^I_k$, c'est-à-dire de la forme $V_S$, est
ce que la page appelle un espace algébrique affine. On a $V_S = V_J$, où $J$ est
l'idéal engendré par $S$, et c'est le système d'équations le plus
naturel.\footnote{En marge, un renvoi au théorème de la base de Hilbert : pour
$k$ noethérien, $k[T_1, \ldots, T_n]$ est noethérien, donc un système fini
suffit.} Avec $A_J = P_I / J$ :
\[
V_J(k') = \{ u : P_I \to k' \mid u(J) = 0 \} \simeq \operatorname{Hom}_{k\text{-alg}}(A_J, k') ,
\]
c'est-à-dire $V_J = V_{A_J}$. Toute $k$-algèbre étant quotient d'une algèbre de
polynômes, les foncteurs $V_S$ sont, à isomorphisme près, exactement les
foncteurs représentables $V_A$ ; se donner un isomorphisme $V_A \simeq V_S$
revient à se donner un isomorphisme $A \simeq P_I / J$. Plus précisément, un
monomorphisme de foncteurs $V_A \to \mathbb{E}^I_k$ correspond à un épimorphisme
d'anneaux $P_I \to A$, qui n'est pas nécessairement surjectif --- ainsi
$k[T] \to k[T, T^{-1}]$ ; les plongements \emph{fermés} correspondent aux
homomorphismes surjectifs, c'est-à-dire aux systèmes de générateurs
$(a_i)_{i \in I}$ de la $k$-algèbre $A$.\footnote{L'exemple
$k[T] \to k[T, T^{-1}]$ est de nous. Le passage est surchargé : une ligne biffée,
puis la phrase récrite au-dessus.}

\textbf{Corollaire (5).} La catégorie des schémas affines sur $k$ est
anti-équivalente à celle des $k$-algèbres.\footnote{La page écrit « algèbres
sur $A$ », pour « sur $k$ » ; la transcription le signale.} Ceux qui se
plongent comme sous-foncteurs fermés dans un $\mathbb{E}^n_k$, $n$ fini,
correspondent aux $k$-algèbres de type fini. Les sous-foncteurs fermés de
$\mathbb{E}^I_k$ correspondent aux idéaux de $P_I$, en renversant l'ordre :
$V_J \subset V_{J'}$ si et seulement si $J' \subset J$, puisque
\[
J = \{ f \in P_I \mid f(x) = 0 \text{ pour toute } k\text{-algèbre } k' \text{ et tout } x \in V_J(k') \} ;
\]
il suffit en effet de prendre $k' = A_J$ et pour $x$ le point universel.

\pagerange{24}{24}
\subsection*{(6)--(8) Restriction aux corps, Nullstellensatz, fonctions (page 24)}

\textbf{(6)} Que se passe-t-il si l'on ne regarde que les points à valeurs dans
des \emph{corps} ? Notons $V'_S$ la restriction de $V_S$ aux corps sur $k$.
Alors $V'_S = V'_J = V'_{\sqrt{J}}$, et $J \mapsto V'_J$ n'est plus injective ;
mais elle l'est sur les idéaux radicaux, car
\[
\sqrt{J} = \{ f \in P_I \mid f(x) = 0 \text{ pour tout corps } k' \text{ sur } k \text{ et tout } x \in V_J(k') \} .
\]
C'est l'énoncé d'algèbre commutative : dans tout anneau, l'intersection des
noyaux des homomorphismes vers des corps --- l'intersection des idéaux premiers
--- est le nilradical. Donc $\sqrt{J} \subset \sqrt{J'}$ si et seulement si
$V'_{J'} \subset V'_J$.\footnote{La page écrit le radical $\widetilde{J}$ ; on
emploie la notation d'aujourd'hui.}

\textbf{(6')} Les mêmes conclusions valent si l'on se restreint à une classe
plus petite de corps, pourvu qu'elle soit assez grande : pour $I$ fini, les
extensions finies des corps résiduels de $k$ suffisent, et si $k$ est un corps
algébriquement clos, $k$ lui-même suffit. C'est le
\emph{Nullstellensatz}, que la marge nomme : dans une algèbre de type fini sur
un corps, tout idéal premier (ou radical) est intersection d'idéaux
maximaux.\footnote{La phrase de la page est entrecoupée d'illisibles (« pourvu
que … finis que $k$ soit assez grand … ») ; on donne l'énoncé que les
exemples entre parenthèses rendent précis. Pour $k$ anneau quelconque et $I$
fini, il résulte du Nullstellensatz appliqué aux fibres de
$\operatorname{Spec} A \to \operatorname{Spec} k$.}

\textbf{(7)} « Importance des … » : le reste est illisible.

\textbf{(8)} Par Yoneda,
\[
A \simeq \operatorname{Hom}_{k\text{-alg}}(k[T], A) \simeq \operatorname{Hom}(V_A, \mathbb{E}^1_k) :
\]
$A$ est \emph{l'anneau des fonctions} sur le schéma affine $V_A$. Si
$A = P_I / J$ et si $f \in A$ se relève en $g \in P_I$, la fonction
$f_{k'} : V_A(k') \to k'$ est induite par $g_{k'} : k'^I \to k'$. La structure
d'anneau de $A$ provient de la structure d'anneau du foncteur
$\mathbb{E}^1_k$, c'est-à-dire de la droite affine vue comme schéma en
anneaux.\footnote{La marge, en biais et surchargée, esquisse la même chose :
$(f+g)(x) = f(x) + g(x)$, $(fg)(x) = f(x)g(x)$.}

\pagerange{26}{27}
\subsection*{(9) Limites projectives dans $\mathrm{Aff}_k$ (pages 26 et 27)}

Les limites projectives dans $\mathrm{Aff}_k$ correspondent aux limites
inductives dans $\mathrm{Alg}_k$. Le point de vue adapté est celui de
$\mathrm{Aff}_k$ comme sous-catégorie pleine des préfaisceaux : le plongement de
Yoneda commute aux limites projectives qui existent, et les limites se
calculent donc point par point.

\begin{itemize}
  \item Les produits correspondent aux produits tensoriels :
  $\prod_i V_{A_i} = V_{\bigotimes_i A_i}$. L'objet final est $e_k = V_k$ ; ses
  sections dans $X$ sont les points de $X$ à valeurs dans $k$.
  \item Les produits fibrés correspondent à $B \otimes_A C$, produit dans
  $\mathrm{Aff}_{k/S}$, somme dans $\mathrm{Alg}_A$. Le changement de base
  $\mathrm{Aff}_{/S} \to \mathrm{Aff}_{/S'}$ est adjoint à droite de l'oubli
  $\mathrm{Aff}_{/S'} \to \mathrm{Aff}_{/S}$ ; en termes de foncteurs, c'est la
  restriction d'un foncteur sur les $k$-algèbres aux $k'$-algèbres.
  \item Produits et produits fibrés donnent toutes les limites. Le noyau d'une
  double flèche $u, v : X \rightrightarrows Y$ est
  $X \times_{Y \times Y} Y$, pris le long de $(u, v)$ et de la diagonale. Du
  côté des algèbres, le conoyau de $u, v : A \rightrightarrows B$ est $B$ divisé
  par l'idéal engendré par les $u(x) - v(x)$, ou encore
  $B \otimes_{A \otimes A} A$, puisque la multiplication
  $A \otimes A \to A$ est surjective de noyau engendré par les
  $x \otimes 1 - 1 \otimes x$.
  \item Les limites projectives filtrantes correspondent aux limites
  inductives filtrantes d'algèbres, qui se calculent sur les ensembles
  sous-jacents.
\end{itemize}

\emph{Exemples, vers les groupes.} Le noyau d'un homomorphisme de groupes
algébriques affines $G \to G'$ est un groupe algébrique affine ; ainsi
$\mathrm{SL}(n)_k = \operatorname{Ker}\bigl(\mathrm{GL}(n)_k \xrightarrow{\det} \mathbb{G}_{m,k}\bigr)$.
Si $G$ opère sur $X$ et $x \in X(k) = \operatorname{Hom}(e_k, X)$, le
\emph{stabilisateur} $G_{(x)}$ est le produit fibré de $G \to X$,
$g \mapsto g.x$, et de $x : e_k \to X$. Pour plusieurs points $x_i \in X_i(k)$,
l'intersection des stabilisateurs --- un produit fibré encore --- est le
stabilisateur de $(x_i)$ dans $\prod X_i$. Ce sont des \emph{schémas en groupes
affines} définis par leur foncteur des points.

\pagerange{27}{28}
\subsection*{(10) Limites inductives dans $\mathrm{Aff}_k$ (pages 27 et 28)}

Les limites inductives dans $\mathrm{Aff}_k$ ne s'interprètent pas commodément
en termes de foncteurs sur les algèbres et n'ont guère de sens géométrique,
sauf dans des cas très particuliers. Elles correspondent aux limites projectives
d'algèbres, qui se calculent sur les ensembles --- « nous fait une belle
jambe ! ».

\emph{L'exemple.} Soient $G = \mathrm{GL}(n)_k$ et $B \subset G$ le sous-groupe
des matrices triangulaires. $B$ définit une relation d'équivalence sur $G$, et
le quotient $G/B$ calculé dans $\mathrm{Aff}_k$ est $e_k$ --- « ce qui est
idiot ! ». En effet ce quotient est le spectre de l'algèbre des fonctions sur
$G$ invariantes par $B$, qui sont les fonctions globales sur la variété des
drapeaux ; celle-ci est projective et connexe, et ses seules fonctions globales
sont les constantes.\footnote{La justification est de nous. L'exemple montre ce
que la page veut montrer : le bon quotient $G/B$ est un schéma projectif, non
affine, et la catégorie affine est trop petite pour les quotients.}

\emph{Sommes.} Les sommes dans $\mathrm{Aff}_k$ correspondent aux produits
d'algèbres. Seules les sommes finies ont une interprétation géométrique
simple. Pour $I$ fini et $X = V_{\prod A_i}$, l'application
\[
\coprod_i \operatorname{Hom}_{k\text{-alg}}(A_i, k') \longrightarrow \operatorname{Hom}_{k\text{-alg}}\Bigl(\prod_i A_i,\ k'\Bigr)
\]
est injective si $k' \neq 0$, mais pas bijective en général : un homomorphisme
$\prod A_i \to k'$ est la donnée d'une famille $(e_i)$ d'idempotents de $k'$
orthogonaux deux à deux et de somme $1$, et d'homomorphismes
$A_i \to k' e_i$.\footnote{La page n'écrit pas « $I$ fini » dans cette phrase,
bien qu'elle vienne de le dire pour l'interprétation géométrique. Pour $I$
infini la description est fausse : pour $k$ fini, les ultrafiltres non principaux
sur $I$ donnent des homomorphismes $k^I \to k$ qui annulent tous les
idempotents $e_i$.} Si $k'$ est \emph{connexe}, c'est-à-dire non nul et sans
idempotent autre que $0$ et $1$ --- par exemple intègre, ou local, ou tel que
$k' / \mathrm{Nil}(k')$ soit intègre ---, l'application est bijective.

Pour $I = \emptyset$ on trouve l'objet initial $\emptyset_k = V_0$, spectre de
l'algèbre nulle : $\emptyset_k(k')$ est vide si $k' \neq 0$ et réduit à un point
si $k' = 0$. C'est un objet initial \emph{strict} --- tout objet qui a une flèche
vers lui lui est isomorphe ---, le « $k$-espace algébrique vide ».

\emph{Conoyaux.} Le conoyau de $u, v : R \rightrightarrows X$ dans
$\mathrm{Aff}_k$ correspond au noyau de $A \rightrightarrows B$ ; il n'a en
général pas beaucoup de sens. La page le juge « O.K. » si $R$ est une relation
d'équivalence sur $X$ finie sur $X$, et renvoie à « la question des passages au
quotient », qui n'est pas traitée dans le dossier.\footnote{Le théorème qui rend
cet « O.K. » précis --- le quotient est effectif et $X \to X/R$ est fini ---
suppose $R$ fini et \emph{localement libre} sur $X$ (SGA 3, exposé V). Sans
platitude, on ne peut pas affirmer davantage que la page.}

Enfin, les sous-schémas fermés d'un schéma affine $V_A$ correspondent aux idéaux
de $A$, l'ordre étant renversé.

\pagerange{28}{29}
\section*{VI. Le spectre (pages 28 à 31, et 21)}

La rubrique (11) associe à un schéma affine $\mathfrak{X} = V_A$ un objet
géométrique, en trois temps : un ensemble, une topologie, un faisceau
d'anneaux.\footnote{Le numéro cerclé est surchargé ; on lit (11).}

\subsection*{(a) L'ensemble sous-jacent (pages 28 et 29)}

Deux points géométriques de $\mathfrak{X}$, à valeurs dans des corps, ont même
« lieu » s'ils proviennent d'un même point à valeurs dans un corps plus petit ;
l'ensemble des lieux est en bijection avec l'ensemble
$\operatorname{Spec} A$ des idéaux premiers de $A$, le point
$u : A \to k'$ ayant pour lieu $\operatorname{Ker} u$. À chaque point $x$ est
attaché son corps résiduel $k(x)$, qui représente le foncteur « ensemble des
points à valeurs dans le corps $k'$ de lieu $x$ » :
ces points sont les $k$-plongements $k(x) \to k'$.

\pagerange{29}{30}
\subsection*{(b) La topologie de Zariski (pages 29 et 30)}

Les fermés de $X = \operatorname{Spec} A$ sont les $V(S)$, pour $S \subset A$ :
$V(S)$ est l'ensemble des idéaux premiers qui contiennent $S$, c'est-à-dire des
lieux des points géométriques où tous les éléments de $S$
s'annulent.\footnote{La page écrit ensuite que $\widetilde{J}_S$ est
l'intersection des idéaux premiers « qui contiennent $A$ » : il faut lire
$S$. La transcription le signale.} On a
$V(S) = V(J_S) = V(\sqrt{J_S})$, où $J_S$ est l'idéal engendré, et il est bon de
se borner aux idéaux radicaux : $V(S) \subset V(S')$ si et seulement si
$\sqrt{J_S} \supset \sqrt{J_{S'}}$, ce qui est l'énoncé (6) de la page 24
relu dans le spectre. De plus
\[
V\Bigl(\bigcup_i S_i\Bigr) = V\Bigl(\sum_i J_i\Bigr) = \bigcap_i V(J_i) \quad (I \text{ quelconque}),
\]
\[
V\Bigl(\bigcap_i J_i\Bigr) = V\Bigl(\prod_i J_i\Bigr) = \bigcup_i V(J_i) \quad (I \text{ fini}),
\]
la seconde parce qu'un premier qui contient un produit fini d'idéaux contient
l'un d'eux. La somme de deux idéaux radicaux n'est pas nécessairement radicale,
note la page ; ainsi $(y) + (y - x^2) = (y, x^2)$ dans $k[x, y]$.\footnote{L'exemple
est de nous.}

C'est donc une topologie. L'adhérence d'un point $x$, d'idéal $\mathfrak{p}_x$,
est $V(\mathfrak{p}_x)$, et $x \in \overline{\{y\}}$ si et seulement si
$\mathfrak{p}_x \supset \mathfrak{p}_y$. Les ouverts
$X_f = X \setminus V(f)$ forment une base stable par intersections finies :
$X_1 = X$, $X_{fg} = X_f \cap X_g$. L'espace $X$ est quasi-compact : si
$A = \sum_{i \in I} J_i$, il existe une partie finie $I_0 \subset I$ telle que
$A = \sum_{i \in I_0} J_i$.\footnote{La page appelle $J$ cette partie finie, la
lettre des idéaux ; on la renomme.} Si $J$ est un idéal,
$\operatorname{Spec}(A/J)$ s'identifie au sous-espace fermé $V(J)$, qui ne
change pas si l'on remplace $J$ par $\sqrt{J}$ : tout fermé est homéomorphe au
spectre d'un anneau réduit.

\textbf{Proposition.} Si $A$ est réduit, $\operatorname{Spec} A$ est
irréductible si et seulement si $A$ est intègre. En effet, $X = V(J') \cup V(J'')$
équivaut à $J'J'' \subset \mathrm{Nil}(A) = 0$, et $V(J') = X$ à $J' = 0$.

\textbf{Corollaire.} Pour $A$ quelconque, les fermés irréductibles de $X$ sont
les $V(\mathfrak{p})$, $\mathfrak{p}$ premier, en bijection décroissante avec
les idéaux premiers, et $V(\mathfrak{p})$ est l'adhérence du point
$\mathfrak{p}$. En particulier tout fermé irréductible a un unique point
générique : $X$ est \emph{sobre}.

Avec la quasi-compacité des $X_f \simeq \operatorname{Spec} A_f$, que la marge de
la page 30 rappelle, ce sont exactement les propriétés qui caractérisent
aujourd'hui les \emph{espaces spectraux}.\footnote{Le nom et la caractérisation
sont de Hochster (1969). Le mot « sobre » est sur la page.} Enfin, l'espace
$\operatorname{Spec} A$ ne dépend que de la structure d'anneau de $A$, et non de
sa structure de $k$-algèbre.

\pagerange{30}{31}
\subsection*{(c) Le faisceau structural (pages 30 et 31)}

« On ne peut reconstituer $A$ à partir du seul espace topologique » : tous les
corps ont le même spectre. Il faut donc un faisceau d'anneaux, et la page le
construit sur la base des $X_f$.

On a $X_f \subset X_g$ si et seulement si $f \in \sqrt{gA}$, c'est-à-dire s'il
existe $n \in \mathbb{N}$ et $h \in A$ avec $f^n = gh$ ; on écrit $f \prec g$.
C'est un préordre. Un préfaisceau sur la base est la donnée d'un ensemble
$P(f)$ pour chaque $f$ et d'une restriction $P(g) \to P(f)$ pour $f \prec g$,
transitives, et l'identité pour $f = g$. On a $X_f = \bigcup_{i \in I} X_{f_i}$
si et seulement si $f_i \prec f$ pour tout $i$ et s'il existe $n$ et des
$h_i$, presque tous nuls, avec $f^n = \sum_i f_i h_i$ ; et le préfaisceau est
un faisceau si, pour tout tel recouvrement, la suite
\[
P(f) \longrightarrow \prod_i P(f_i) \rightrightarrows \prod_{i,j} P(f_i f_j)
\]
est exacte.\footnote{La page écrit « $f_i \in I$ » pour « $i \in I$ » ; la
transcription le signale. Elle note ensuite qu'on peut supposer $I$ fini, par
quasi-compacité.}

Soit $M$ un $A$-module. Pour $f \in A$, on pose
$M_f = \varinjlim_n M$, les flèches de transition étant la multiplication par
$f^{n'-n}$ ; les éléments de $M_f$ sont les $m/f^n$, avec les règles de calcul
habituelles, et $M \to M_f$ est la solution universelle du problème de rendre
$f$ inversible. On a $M_f \simeq M \otimes_A A_f$.

\textbf{Lemme.} Si $f \prec g$, il existe un unique homomorphisme
$M_g \to M_f$ compatible aux applications canoniques ; car $f^n = gh$ étant
inversible sur $M_f$, $g$ l'est aussi.

On obtient ainsi un préfaisceau $\widetilde{M}$ sur la base des $X_f$, qui est un
préfaisceau de modules sur le préfaisceau d'anneaux $f \mapsto A_f$. C'est un
\emph{faisceau} : la suite
$M_f \to \prod_i M_{f_i} \rightrightarrows \prod_{i,j} M_{f_i f_j}$ est exacte
dès que $X_f = \bigcup X_{f_i}$. Comme $(M_f)_g \simeq M_g$ pour $g \prec f$, on
se ramène au cas $f = 1$, $\sum f_i h_i = 1$, qui est l'argument classique de
partition de l'unité.\footnote{La page s'arrête, après la réduction au cas
$f = 1$, sur des points de suspension ; la démonstration n'est pas écrite. Le mot
qui précède l'énoncé (« Prop » ou « Th ») est illisible.}

En particulier, $X$ est annelé par $\mathcal{O}_X = \widetilde{A}$, faisceau de
$A$-algèbres. Sa fibre en $x$, d'idéal $\mathfrak{p}$, est la limite inductive
des $A_f$ pour $f \notin \mathfrak{p}$ ; comme $A \setminus \mathfrak{p}$ est
filtrant pour $\succ$, c'est le localisé $A_{\mathfrak{p}}$, et c'est un anneau
local. De même la fibre de $\widetilde{M}$ est $M_{\mathfrak{p}}$.

\pagerange{21}{21}
\subsection*{Fonctorialité et topologie des spectres (page 21)}

La page 21, sous le titre « Topologie des spectres », appartient à cette
rubrique, et plus précisément à la « 10') Fonctorialité de $\operatorname{Spec} A$ »
qu'annonce la marge de la page 30 : elle emploie les lettres gothiques et les
$V_A$ du plan, et suppose connue la base des $X_f$.

Un espace est \emph{noethérien} si toute suite décroissante de fermés est
stationnaire ; il est alors quasi-compact, et tout fermé n'a qu'un nombre fini
de composantes irréductibles.\footnote{La page définit l'espace noethérien comme
« quasicompact tel que » ces deux conditions ; la seconde entraîne les deux
autres, et la définition est donc équivalente à l'usuelle.} Dans un espace dont
les composantes irréductibles sont en nombre fini --- ou plus généralement
forment une famille localement finie ---, les composantes connexes sont
ouvertes.

Soient $u : A \to B$ et $\varphi : X = \operatorname{Spec} B \to Y = \operatorname{Spec} A$.
Alors\footnote{La page écrit $\varphi(x_{\mathfrak{p}}) = y_{\varphi^{-1}(\mathfrak{p})}$ :
l'image réciproque de l'idéal est prise par $u$. Les indices des deux lignes
suivantes sont surchargés, $\varphi(f)$ corrigé en $u(f)$.}
\[
\varphi(\mathfrak{q}) = u^{-1}(\mathfrak{q}), \qquad
\varphi^{-1}(V(J)) = V(u(J)B), \qquad
\varphi^{-1}(Y_f) = X_{u(f)} .
\]

\begin{itemize}
  \item[a)] Si $u$ est surjectif, $\varphi$ est une immersion fermée.
  \item[b)] Pour $u : A \to A_f$, $\varphi$ est une immersion ouverte, et
  $\operatorname{Spec}(A_f) \simeq Y_f$. Dans ces deux cas les extensions
  résiduelles sont triviales ; il en va de même dès que $u$ est un épimorphisme
  d'anneaux, c'est-à-dire dès que $V_B \to V_A$ est un monomorphisme.
  \item[c)] Pour $I$ \emph{fini},
  $\operatorname{Spec} \prod_{i \in I} A_i \simeq \coprod_i \operatorname{Spec} A_i$.\footnote{La
  page écrit $\prod_I$ sans condition. Pour $I$ infini, le spectre du produit est
  beaucoup plus gros que la somme : il contient un point pour chaque
  ultrafiltre sur $I$.}
  \item[d)] Le foncteur « espace sous-jacent » ne commute pas aux produits
  fibrés : l'application
  \[
  \lvert \mathfrak{X} \times_S \mathfrak{Y} \rvert \longrightarrow
  \lvert \mathfrak{X} \rvert \times_{\lvert S \rvert} \lvert \mathfrak{Y} \rvert
  \]
  est surjective mais non injective ; sa fibre au-dessus de $(x, y)$, avec $x$ et
  $y$ au-dessus de $s$, est $\operatorname{Spec}\bigl(k(x) \otimes_{k(s)} k(y)\bigr)$,
  qui n'est pas vide. Ses points correspondent aux corps composés $K'$ de $k(x)$
  et $k(y)$ au-dessus de $k(s)$ ; la page les dessine en quatre croquis, et la
  marge renvoie à l'exemple du plan $\mathbb{E}^1 \times \mathbb{E}^1$.
\end{itemize}

Sur un corps algébriquement clos et pour des schémas de type fini, c'est le
\emph{spectre maximal} qui commute aux produits finis, en tant qu'ensemble ---
ses points sont les points rationnels ---, mais sa topologie n'est pas la
topologie produit.\footnote{La page porte, à la fin de d), « [vrai sur un corps
alg. clos] » ; la phrase suivante précise qu'il s'agit du spectre maximal, et
c'est ainsi qu'on la lit. Pour le spectre tout entier l'énoncé serait faux même
sur un corps algébriquement clos : le point générique du plan et celui de la
diagonale ont la même image.}

\pagerange{19}{20}
\section*{VII. Compléments : sommes, limites et spectres dans $\mathrm{Aff}_k$ (pages 19 et 20)}

La feuille est titrée de sa main « Géom. Algébrique » et annonce des
« compléments sur le formalisme des $\varprojlim$ et des sommes dans
$\mathrm{Aff}_k$ ». Elle complète les rubriques (9) et (10), dont elle emploie les
notations, et suppose le spectre de la rubrique (11).

\emph{a) Les sommes finies sont disjointes et universelles.} Soit
$X = \coprod_{i \in I} X_i$, $I$ fini, c'est-à-dire $A = \prod_i A_i$. Alors :
chaque $X_i \to X$ est un monomorphisme, et même une immersion fermée, puisque
$A \to A_i$ est surjectif ; $X_i \times_X X_j = \emptyset$ pour $i \neq j$,
puisque $A_i \otimes_A A_j = 0$ (l'idempotent de $A_i$ vaut $1$ dans $A_i$ et
$0$ dans $A_j$) ; et pour tout $X' \to X$, $X' \simeq \coprod_i X'_i$ avec
$X'_i = X' \times_X X_i$, puisque toute $A$-algèbre $B$ se décompose en
$\prod_i B e_i$. Donc les produits, et les produits fibrés, sont distributifs
par rapport aux sommes finies.

Il en résulte que le foncteur
\[
(\text{ensembles finis}) \longrightarrow \mathrm{Aff}_k, \qquad
\Gamma \longmapsto \Gamma_k = \coprod_{\Gamma} e_k = V_{k^{\Gamma}}
\]
commute aux produits finis et aux produits fibrés, donc aux limites projectives
finies. Il transforme donc un monoïde (un groupe) fini en un monoïde (un groupe)
algébrique affine : c'est le \emph{schéma en groupes constant} $\Gamma_k$. Son
algèbre est l'algèbre $k^{\Gamma}$ des fonctions sur $\Gamma$, contravariante
en $\Gamma$ ; si $\Gamma$ a une loi $\Gamma \times \Gamma \to \Gamma$, elle
définit la comultiplication
$k^{\Gamma} \to k^{\Gamma \times \Gamma} \simeq k^{\Gamma} \otimes k^{\Gamma}$.\footnote{La
page note aussi cette algèbre $k[I]$, en prévenant contre la confusion avec les
polynômes. On l'écrit $k^{\Gamma}$ : aujourd'hui $k[\Gamma]$ désigne d'ordinaire
l'algèbre \emph{du groupe}, dont la structure d'algèbre de Hopf est duale de
celle-ci ; pour $\Gamma$ commutatif, son spectre est le groupe diagonalisable
dual de $\Gamma_k$, et non $\Gamma_k$. La page emploie tantôt $\Gamma$, tantôt
$I$, pour le même ensemble.} Le foncteur $\Gamma_k$ associe à $k'$ l'ensemble
des décompositions de $1$ en idempotents orthogonaux indexées par
$\Gamma$,\footnote{C'est la marge ; le mot qui qualifie ces décompositions est
lu « parfaites » avec doute : il s'agit des systèmes \emph{complets}, de somme
$1$.} ce qui redonne l'ensemble $\Gamma$ lorsque $k'$ est connexe ; et
$\operatorname{Spec} k^{\Gamma} = \coprod_{\Gamma} \operatorname{Spec} k$.

\emph{b) Produits fibrés le long d'un monomorphisme.} Si l'un des deux
morphismes est un monomorphisme, le carré des espaces sous-jacents est encore
cartésien, en tant que carré d'ensembles : les extensions résiduelles d'un
monomorphisme sont triviales (page 21), et les fibres
$\operatorname{Spec}\bigl(k(x) \otimes_{k(y)} k(y')\bigr)$ sont réduites à un
point. Si de plus $\mathfrak{Y}' \to \mathfrak{Y}$ est une immersion fermée, une
immersion ouverte, ou plus généralement correspond à une localisation
$A \to S^{-1}A$ --- trois classes stables par changement de base, et formées
d'homéomorphismes sur l'image ---, la topologie de $X'$ est la topologie induite
par celle de $X$.\footnote{La page dit « plus généralement si celle correspond à
$A \to A S^{-1}$, disons l'homéomorphisme dans ». On s'en tient aux trois cas
qu'elle nomme, pour lesquels l'énoncé est vrai ; pour un monomorphisme qui n'est
qu'un homéomorphisme sur son image, la stabilité par changement de base ne va
pas de soi. La page dessine deux carrés : à gauche celui des espaces
algébriques affines, en lettres gothiques $\mathfrak{X}$, $\mathfrak{Y}$ ; à
droite celui des spectres, en romaines $X$, $Y$ ; et la phrase emploie les
mêmes lettres, $\mathfrak{Y}' \to \mathfrak{Y}$ pour le morphisme et $X'$ pour
l'espace. On les garde. Les flèches du croquis de gauche sont en pointillé, et
celle du bas porte deux pointes.} Cas particulier : l'intersection de deux sous-objets.

\emph{c) Adhérence de l'image.} Pour $u : A \to B$ et
$\varphi : X = \operatorname{Spec} B \to Y = \operatorname{Spec} A$, et
$Z = V(J) \subset X$, on a $\overline{\varphi(Z)} = V(u^{-1}(J))$. En
particulier $\varphi$ est dominant si et seulement si
$\operatorname{Ker} u \subset \mathrm{Nil}(A)$.\footnote{La page écrit
$\operatorname{Ker} \varphi$ ; il s'agit du noyau de $u$. La transcription le
signale. La formule pour $\overline{\varphi(Z)}$ n'est qu'annoncée sur la page ;
on la donne.}

\pagerange{15}{17}
\section*{VIII. Feuillets étrangers au cours : le produit de deux topos (pages 17 et 15), et la page 25}

Les pages 15 et 17 ne parlent pas de géométrie algébrique. Elles portent sur une
question de théorie des topos, et on les lit dans l'ordre de l'argument : la page
17, qui reprend en cours de raisonnement, puis la page 15, qui la
continue.\footnote{La page 15 est un fragment isolé en haut d'une feuille, qui
commence par « Or » et emploie le $C'$ et le $\widetilde{D}$ introduits à la
page 17 ; elle en utilise aussi le corollaire (que $\gamma$ est un plongement).
Les deux feuilles sont traversées de longs traits au crayon.}

Le cadre est posé ailleurs. Dans le dossier 161-6, à la page 22 --- une page
barrée d'une croix, mais lisible ---, on lit : pour un topos $E$ et une petite
catégorie $C$, $\pi(C, E) = \underline{\operatorname{Hom}}(C^{\circ}, E)$ est un
topos, qui dépend 2-fonctoriellement de $(C, E)$ ; on définit un morphisme de
topos $\gamma(C, E) : \pi(C, E) \to \widehat{C} \times E$ vers le produit dans
la 2-catégorie des topos ; et l'on demande si $\gamma(C, E)$ est toujours une
équivalence, en notant trois cas : les sommes, les groupoïdes, et
$E = \widehat{D}$, où $\pi(C, E) \simeq \widehat{C \times D}$.\footnote{Dossier
161-6, page 22 (transcription \texttt{161-6/batch-02.fr.tex}). On rapproche les
deux dossiers par leur contenu, qui est le même problème dans les mêmes
notations ; on ne les fusionne pas, et rien ne dit que les feuilles aient été
écrites au même moment.} La page 17 de notre dossier s'y enchaîne par le
contenu : elle s'ouvre sur « donc la question est si c'est une équiv. de
topos », traite le cas $E = \widehat{D}$, puis le cas général. Rien sur les
feuillets ne dit laquelle des deux pages a été écrite la première.\footnote{Un
détail ne s'accorde pas avec une suite directe : la page 17 désigne le cas
général par « c) », la lettre que la page 22 du dossier 161-6 donne au cas
$E = \widehat{D}$ (ses cas sont a) les sommes, b) les groupoïdes,
c) $E = \widehat{D}$). La page 17 numérote aussi sa question « (2) », après le
« (1) » encadré de 161-6 ; les deux indices tirent en sens contraires, et on
ne tranche pas.}

\pagerange{17}{17}
\subsection*{Le cas des préfaisceaux, puis la réduction (page 17)}

Pour $E = \widehat{D}$, la question devient :
\[
(2) \qquad \widehat{C \times D} \overset{?}{\simeq} \widehat{C} \times \widehat{D} .
\]
Si $C$ et $D$ ont des limites projectives finies, $C \times D$ aussi, et l'on
conclut par la propriété universelle des topos de préfaisceaux : pour une
catégorie $A$ à limites finies et un topos $E$,
\[
\underline{\operatorname{Homtop}}(E, \widehat{A}) \simeq
\bigl(\underline{\operatorname{Hom}}_{\mathrm{lex}}(A, E)\bigr)^{\circ},
\qquad u \longmapsto u^{*}|_A ,
\]
les morphismes de topos $E \to \widehat{A}$ correspondant aux foncteurs exacts
à gauche $A \to E$.\footnote{C'est le théorème de Diaconescu dans le cas où $A$
a des limites finies ; la page le rappelle sans nom. Le signe ${}^{\circ}$ tient à
la convention choisie pour les 2-morphismes de topos, par les images directes ou
par les images inverses ; il est sans effet sur la question posée.} Un foncteur
exact à gauche sur $C \times D$ est la même chose qu'un couple de foncteurs
exacts à gauche sur $C$ et sur $D$, et (2) est une équivalence.

Dans le cas général, on choisit un plongement pleinement fidèle $i : C \to C'$
dans une petite catégorie à limites finies, et un plongement de topos
$j : E \hookrightarrow \widehat{D}$, c'est-à-dire une présentation
$E \simeq \widetilde{D}$ par un site $D$, qu'on peut prendre à limites
finies. On obtient un carré commutatif à isomorphisme près :

\begin{tikzcd}
\pi(C, \widetilde{D}) \arrow[r, "{\alpha = \pi(i, j)}"] \arrow[d, "\gamma"'] & \pi(C', \widehat{D}) \simeq \widehat{C' \times D} \arrow[d, "\gamma'"] \\
\widehat{C} \times \widetilde{D} \arrow[r, hook, "{\beta = i \times j}"'] & \widehat{C'} \times \widehat{D}
\end{tikzcd}

où $\gamma'$ est une équivalence par le cas précédent, et $\beta$ un plongement
de topos, produit de deux plongements. Le morphisme $\alpha$ se factorise en
\[
\pi(C, \widetilde{D}) \xrightarrow{\ \alpha_0\ } \pi(C, \widehat{D}) \simeq \widehat{C \times D}
\xrightarrow{\ \alpha_1\ } \pi(C', \widehat{D}) \simeq \widehat{C' \times D} ;
\]
$\alpha_1$ est un plongement parce que $C \times D \to C' \times D$ est
pleinement fidèle, et $\alpha_0$ en est un parce que son image directe,
$u \mapsto j_* \circ u$, est pleinement fidèle.

\textbf{Corollaire.} $\gamma(C, E) : \pi(C, E) \to \widehat{C} \times E$ est un
plongement de topos. En effet $\beta\gamma \simeq \gamma'\alpha$ est un
plongement, et $\beta$ aussi, donc l'image directe de $\gamma$ est pleinement
fidèle.

Reste à montrer que $\gamma$ est essentiellement surjectif, c'est-à-dire que
l'image de $\pi(C, \widetilde{D})$ dans $\widehat{C' \times D}$ est tout le
sous-topos $\widehat{C} \times \widetilde{D}$. La page propose de décrire ce
sous-topos comme la sous-catégorie pleine des $F$ tels que $p_*(F)$ soit dans
l'image essentielle de $\widehat{C}$ et $q_*(F) \in \widetilde{D}$, où $p$ et $q$
sont les deux projections. Cette description est trop faible. La bonne est :
$F(c, -) \in \widetilde{D}$ pour tout objet $c$ de $C'$, et $F(-, d)$ dans
l'image essentielle de $\widehat{C}$ pour tout objet $d$ de $D$.\footnote{Contre-exemple
à la description de la page, dans ses propres hypothèses : $C = C'$ l'ensemble
ordonné $\{0 < 1\}$, $D$ la catégorie ponctuelle, et $\widetilde{D}$ le topos
dégénéré (la topologie pour laquelle le crible vide couvre), dont le seul objet
est l'objet final. Alors $\widehat{C} \times \widetilde{D}$ est dégénéré ; mais
$q_*(F) = F(1)$, puisque $1$ est final dans $C$, et la condition de la page ne
demande que $F(1)$ réduit à un point, laissant $F(0)$ arbitraire.} C'est
exactement l'image de $\pi(C, \widetilde{D})$ --- ce que la page voulait
établir, « ce qui prouvera que $\gamma$ est une équivalence ». Elle s'arrête sur
ces mots.

\pagerange{15}{15}
\subsection*{Une autre route (page 15)}

La page 15 prend la route correcte. Elle vérifie « directement » que, comme
sous-catégories strictement pleines de $\pi(C', \widehat{D})$,
\[
\pi(C, \widetilde{D}) = \pi(C', \widetilde{D}) \cap \pi(C, \widehat{D}),
\]
et, de l'autre côté, que le produit se décompose de même comme intersection de
sous-topos :
\[
\widehat{C} \times \widetilde{D} = \bigl(\widehat{C'} \times \widetilde{D}\bigr) \cap \bigl(\widehat{C} \times \widehat{D}\bigr).
\]
Les deux égalités sont justes : l'intersection de deux sous-topos est la
sous-catégorie des objets qui sont des faisceaux pour les deux topologies, et un
produit de sous-topos est l'intersection de leurs images réciproques. On est
ainsi ramené à montrer que $\gamma(C, \widehat{D})$ et
$\gamma(C', \widetilde{D})$ sont des équivalences, c'est-à-dire, puisque ce sont
des plongements, que leurs images directes sont essentiellement surjectives. La
page s'arrête là.\footnote{Le losange de la page est tracé sans pointes, sauf un
trait. Au-dessus du premier produit « $\times_{\mathrm{Top}}$ », un petit signe
illisible surmonté d'un tilde.}

\emph{La réponse.} La question posée sur ces feuilles et dans le dossier 161-6
a une réponse positive : pour toute petite catégorie $C$ et tout topos $E$, le
topos $\underline{\operatorname{Hom}}(C^{\circ}, E)$ des préfaisceaux sur $C$ à
valeurs dans $E$ est le produit $\widehat{C} \times E$ dans la 2-catégorie des
topos. C'est un fait classique de la théorie des topos --- le changement de base
d'un topos de préfaisceaux est un topos de préfaisceaux ---, qu'on trouve dans les
traités de Johnstone.\footnote{Ce paragraphe est de nous. Il ne dit rien de ce
que Grothendieck savait ou ne savait pas au moment d'écrire ces pages.}

\pagerange{25}{25}
\subsection*{La page 25}

Une feuille au crayon, barrée de longs traits diagonaux croisés, sur un autre
sujet que les pages qui l'entourent : on y lit, par fragments, un topos $X$ et un
objet $F$, le topos induit $X/F$, un
$\underline{\operatorname{Hom}}_X(F, \mathbb{L}_X)$, des puissances $I^E$ indexées
par un ensemble ordonné $E$, et la ligne « Topos associé à l'espace compact
$I^E$ ». Les fragments ne suffisent pas à reconstituer un énoncé, et l'on n'en
propose pas.\footnote{La lettre notée $\mathbb{L}_X$ peut être un $\mathcal{L}$
cursif ; la transcription ne tranche pas. Rien n'indique que cette page ait un
rapport avec les pages 15 et 17, sinon qu'elle parle aussi de topos.}

\end{document}
